\documentclass[english,10pt,a4paper,titlepage]{scrartcl}
\usepackage{report-style}
\setcounter{secnumdepth}{3}
\setcounter{tocdepth}{2}
\title{Rocket GNC Monitor v0}
\subtitle{Complete functionality and developer documentation\\Application version 0.1.0}
\author{MIT Rocket Team\\Avionics and GNC}
\date{21 September 2026}
\begin{document}
\maketitle
\section*{Document control and reading guide}
\addcontentsline{toc}{section}{Document control and reading guide}
This report documents the implemented desktop application in
\code{ui_work/rocket_UI_v0}, including virtual pointer rehearsal, global Settings, bundled OpenRocket-MIT v6.2, Digital/Analog USB video, the portable flight-file workflow and
URRG launch-site preset. Part I explains the operator interface and its
observable behavior. Part II describes the software architecture, contracts,
source files, build process and extension points. Appendices provide configuration,
sample and file-field references.

\begin{longtable}{@{}L{.27\linewidth}L{.68\linewidth}@{}}
\toprule Item & Baseline \\\midrule\endhead
Documentation date & 21 September 2026. \\
Software identity & Rocket GNC Monitor, Python package version 0.1.0, application generation v0. \\
Implementation reviewed & Source files and checked-in resources under \code{rocket_UI_v0}; exact file hashes are recorded in \code{evidence/source-manifest.json}. \\
Current distribution & macOS arm64 app and DMG, rebuilt for smooth antenna/3D flight rendering, MGRS/relative antenna entry, virtual pointer rehearsal, dropdown sizing, Settings, OpenRocket-MIT v6.2, Digital/Analog video, flight save/load and URRG. Windows makefiles remain available; the Windows ZIP predates these changes. \\
Validation evidence & 113-test Mac regression run for virtual pointer, Settings, engine selection and dual video and packaged startup/demo/flight/OpenRocket verification. Historical September 20 results are retained separately. Physical dual-camera acceptance remains outstanding. \\
Report style & KOMA-Script, 10 pt Latin Modern, A4 with 18 mm margins, blue headings, shaded tables and code blocks, matching the supplied GenMAPS report. All diagrams are editable TikZ with LaTeX text. \\
\bottomrule
\end{longtable}
\statusnote{Reading the evidence}{Implemented behavior is distinguished from nominal simulation assumptions and future controlled-vehicle work. A valid software test result is not a measured mechanical pose, a radio acknowledgment or field acceptance.}
\begin{figure}[htbp]\centering
\begin{tikzpicture}[node distance=12mm and 12mm]
\node[box] (rocket) at (0,0) {Rocket flight computer\\Zephyrus telemetry};
\node[box] (ground) at (5.1,0) {Ground station board\\radio and USB serial};
\node[accentbox] (app) at (5.1,-2.05) {Rocket GNC Monitor\\views, commands\\and recording};
\node[box] (pointer) at (10.2,-2.05) {Pointer control board\\separate USB serial};
\node[box] (camera) at (0,-2.05) {USB video receiver\\camera input};
\node[smallbox] (files) at (.2,-4.25) {Flight files\\and session replay};
\node[smallbox] (engine) at (5.1,-4.25) {Private OpenRocket\\nominal simulation};
\node[smallbox] (wind) at (10.2,-4.25) {Wind layers\\manual or online};
\draw[flow] (rocket) -- node[above,font=\footnotesize]{radio} (ground);
\draw[flow] (ground) -- node[right,font=\footnotesize,align=left]{telemetry /\\uplink} (app);
\draw[flow] (camera) -- node[above,font=\footnotesize]{video} (app);
\draw[flow] (app) -- node[above,font=\footnotesize]{aim} (pointer);
\draw[flow] (app.south west) -- (files.north east);
\draw[flow] (engine) -- (app);
\draw[flow] (wind) -- (engine);
\draw[optional] (files.north) to[bend left=18] (app.west);
\end{tikzpicture}
\caption{System boundary. The two serial boards are independent. Simulation, replay and drawing operate inside the desktop application; the app does not infer unavailable actuator feedback. Dashed flow denotes offline restoration.}
\label{fig:system}
\end{figure}

\clearpage
{\small\RedeclareSectionCommand[tocbeforeskip=.35em]{section}\tableofcontents}
\clearpage
\part{Application functionality and operation}
\section{Purpose, scope and operating model}\label{sec:scope}
Rocket GNC Monitor is a native desktop ground application for an amateur rocket.
It combines two independent serial links, Digital and Analog USB video, recorded-flight playback,
a visual antenna assembly, telemetry and actuator displays, and a nominal
OpenRocket trajectory. The current transport contract is Zephyrus. The interface
has places for a configurable number of canards and four fin tabs, but the
Zephyrus packet contains four legacy servo drives; it does not provide the new
vehicle's requested and measured canard/tab angles.

The application starts in \ui{LIVE}, on \ui{Control panel}, without opening either
serial port, starting polling or opening a camera. Connect each board explicitly.
A ground-board USB connection, a working rocket radio link and a pointer-board
USB connection are three separate observations. A screen showing telemetry does
not imply that an uplink command was acknowledged by the rocket.

\begin{longtable}{@{}L{.13\linewidth}L{.38\linewidth}L{.42\linewidth}@{}}
\toprule Mode & Source and controls & Persistence and transition behavior \\\midrule\endhead
LIVE & Physical Zephyrus telemetry; ground commands require its board connection. Pointer manual controls use an independent physical or virtual connection. & Logging captures received data. Source changes close physical transports and clear displayed history. Reconnect and restart polling explicitly. \\
DEMO & One complete bundled GS1, GS2 or GS3 receiver log; synthetic video test pattern. Commands are simulated. & No physical serial transport. Can record a rehearsal or export the complete selected dataset. Returning to LIVE restores the prior reference and retains mission settings. \\
REPLAY & Indexed session telemetry and recorded video, initially paused. Physical commands disabled. & Seek, step, change speed and export. Session files are read for replay; a new flight archive can preserve the session and displayed position. \\
\bottomrule
\end{longtable}

The three source modes are not launch states. Rocket states such as Preflight,
Flight and Main come from telemetry and are shown within a mode. Likewise,
\ui{Start Logging} is independent of \ui{Start Polling}: opening a recording does
not by itself create telemetry.

\subsection{Supported machines and offline use}
The current verified build is macOS arm64. The distributable app includes Python,
Qt/PySide6, NumPy, serial support, FFmpeg, the private Java runtime, the team's
OpenRocket engine and bridge, motors, demo recordings and Lucida Grande. A
production machine does not need a separate Python or Java installation.
The macOS app is packaged with a minimum-system declaration of macOS 13; that is
not evidence that every such system has been tested.

Windows has native build entry points and a previously tested x64 package.
The current Mac feature revision has not been rebuilt and qualified on Windows.
Linux-related camera and runtime branches exist, but Linux is not a qualified
release target. Keep the whole Windows package together; its executable alone is
insufficient.

Serial reception, local video input, demonstration, recording/replay, coordinate
conversion and simulation work offline when the package and inputs are present.
Online wind retrieval requires internet access. USB devices may require an
appropriate OS driver, and macOS camera permission must be granted.

\subsection{What the app establishes}
A command marked \emph{sent} has been dispatched on the serial interface; neither
legacy protocol provides the command-level acknowledgment needed to prove the
requested physical action. The antenna graphic follows the last sent target,
not measured motor encoders. Its dimensions are illustrative rather than
surveyed CAD.

The OpenRocket output is a nominal reference. It does not qualify the new
controlled rocket's aerodynamic, actuator or control-loop behavior. Integrated
Zephyrus rotations are displayed as received and are not a validated attitude
solution. Missing measurements are shown as unavailable rather than inferred.

\section{A complete operating workflow}\label{sec:workflow}
\subsection{Before connecting equipment}
\begin{enumerate}
\item Start the app and confirm \ui{LIVE} on the connection panel. Both board
cards should be disconnected.
\item Open \ui{Mission $\rightarrow$ Configure}. Name the mission and enter the
launch origin, or choose \ui{URRG}. Enter the ellipsoid and mean-sea-level
elevations separately, then select the verified legacy height interpretation.
\item In \ui{Mission \& wind}, enter and apply a manual wind profile or retrieve
one for the intended UTC time. Select the rocket's \code{.ork} file and any
additional motor curves. Run the nominal simulation if a reference is needed.
\item Save a planning \code{.rktflight} file so the settings, selected model/motors
and available reference can be reopened together.
\end{enumerate}
\subsection{Receiving and recording}
\begin{enumerate}
\item Refresh the two serial selectors. Assign the ground-station board and
antenna-pointer board to different ports.
\item Connect the ground board and click \ui{Start Polling}. Wait for
\ui{Rocket telemetry live}, then verify packet age and the displayed data.
\item Connect the pointer board if pointing control is needed. Manual pointer
control is independent of the ground board.
\item On \ui{Flight overview}, choose \ui{Find cameras}, select the USB receiver
and press \ui{Start camera}. Confirm that frames are arriving.
\item Click \ui{Start Logging}. It immediately creates a unique session folder.
Check Diagnostics for the recording path and any loss/error indication.
\item For tracking, use \ui{Send to AntPtr} to freeze the receiver GPS, then
\ui{Track live rocket}. Tracking requires both boards, polling and a usable
fresh target. Review the meaning of ZERO and Stop tracking in
Section~\ref{sec:antenna}.
\end{enumerate}
\subsection{After the flight}
Stop tracking and stop logging. Use \ui{Save flight} to create the portable
archive; the save waits for pending recorder/video closure. Reopen the file to
check the recorded session in paused replay. Retain the original session folder
when preserving full capture provenance. Close the app normally to finalize its
workers. A flight saved while logging is still active is a snapshot and can omit
the unfinished video segment.

\section{Global settings}\label{sec:settings}
Open \ui{Rocket GNC Monitor > Settings} on Mac or press Command-comma.
Windows uses \ui{File > Settings} or Ctrl-comma. The window is available in every
workspace and source mode, including disconnected LIVE. Reopening it brings the
existing editor forward.

\begin{longtable}{@{}L{.24\linewidth}L{.69\linewidth}@{}}
\toprule Preference & Behavior \\\midrule\endhead
OpenRocket JAR & Browse to a complete compatible OpenRocket-MIT JAR. \ui{Use bundled} clears the override and shows the bundled release and path. Java remains included with the app. \\
Simulation time limit & 10--1800 s; default 120 s. Engine and time-limit changes apply to the next job; the running job retains its snapshot. \\
Recording folder & Destination for new logging sessions. \ui{Use default} restores the application's sessions directory. Existing files and active recordings retain their locations. \\
Appearance & Dark or Daylight. The \ui{View > Daylight theme} action updates the same persistent preference. \\
\bottomrule
\end{longtable}
\ui{Save} validates and stores changes. \ui{Cancel} discards edits.
\ui{Restore Defaults} changes the form and requires Save to take effect.
A missing or unreadable custom JAR is reported; select the relocated file or use
the bundled engine. The current Mac package includes OpenRocket-MIT v6.2 and
requires no separate OpenRocket installation or runtime download.

Preferences are stored in \code{settings.json} under the per-user application
data directory; the exact path is shown in the window. \code{--data-dir} also
selects a separate settings store. These machine-local paths and preferences
are not exported in mission JSON or portable flight files. Site, wind, rail,
alert thresholds and video offset remain mission-specific. Startup remains
LIVE with hardware closed regardless of the saved preferences.

\section{Navigation and connection control}
The persistent header contains the mission name, source-mode selector,
Start/Stop Polling and Start/Stop Logging. In LIVE and REPLAY the two serial cards
remain visible; controls are enabled according to mode and board state. The
status area provides local wall-clock time, source/status messages and alerts.
The \ui{View} menu selects the daylight theme; this changes contrast, not data.

\begin{longtable}{@{}L{.23\linewidth}L{.70\linewidth}@{}}
\toprule Workspace & Function \\\midrule\endhead
Control panel & Ground USB, rocket radio and pointer USB status; packet age, RSSI, valid/rejected counts and readiness. \\
Flight overview & Main metrics, camera/file video, projected trajectory, vertical profile and attitude instrument. \\
Antenna pointer & Orbitable assembly, manual azimuth/elevation, direction commands, ZERO, ground-GPS freeze and tracking. \\
GNC \& actuators & Three angular-rate and integrated-rotation traces, canard/tab channel slots and legacy servo drives. \\
Mission \& wind & Mission configuration, launch location, wind layers, weather, rocket/motor selection and simulation. \\
Sessions \& replay & Portable flight files, session-folder replay, legacy CSV import, sample export, seeking and event timeline. \\
Diagnostics & Decoded sample JSON, packet counters, recording state, events and alert acknowledgment. \\
Rocket controls & Four subviews: Telemetry, Rocket commands, Power, Recovery \& pyros. \\
\bottomrule
\end{longtable}

\subsection{Rocket radio status}
\begin{longtable}{@{}L{.20\linewidth}L{.38\linewidth}L{.35\linewidth}@{}}
\toprule Status & Meaning & Normal next action \\\midrule\endhead
Disconnected & Ground USB transport is not connected. & Select and connect the telemetry board. \\
Paused & Ground USB is open but telemetry polling is stopped. & Start polling to receive; uplinks remain available. \\
Waiting & Polling is active but no fresh valid packet has established reception. & Check transmitter, receiver, antenna and serial source. \\
Receiving & A checksum-valid LIVE packet arrived within the configured freshness interval. & Monitor packet age and loss indicators. \\
Stale & Previously received valid telemetry has aged beyond that interval. & Restore reception; explicitly restart tracking if it stopped. \\
\bottomrule
\end{longtable}
The default freshness interval is 2 s. Reconnecting the ground board clears prior
reception evidence. A connected USB device alone never proves a rocket radio
link. The trailer carrying ground-receiver GPS has no independent checksum.

\subsection{Serial behavior and command availability}
Both boards use 115200 baud, eight data bits, no parity, one stop bit and no
flow-control override. DTR/RTS follow the original pyserial defaults.
Ground reads use a 1 s timeout; pointer reads use 0.1 s. The selectors reject
assigning the same serial device to both roles, including the usual macOS
\code{/dev/tty.} versus \code{/dev/cu.} alias. A connected or connecting
board's port is removed from the other selector and restored when disconnected.

Stopping polling leaves the port open and leaves rocket uplinks available.
Disconnecting a board closes its transport and clears relevant pending state.
There is no automatic reconnection, retransmission or tracking restart. A
partially written packet is treated as uncertain, not retried.

\section{Flight overview, coordinates and time}\label{sec:flightview}
\subsection{Metrics and plots}
The six main metrics are altitude in metres, velocity in metres per second,
aligned flight time in seconds, battery in volts, receiver RSSI in dBm and sample
age in seconds. Missing values remain blank/unavailable. Battery is the sum of
the three decoded cell voltages; velocity is the legacy reported integrated
velocity rather than a recomputed derivative of the plotted altitude.

The trajectory view offers \ui{Plan: East/North}, \ui{Side: East/Up} and a fixed
projected \ui{Perspective}. The actual trace uses accepted local positions. The
reference is dashed, with a diamond at the simulated position for the same
aligned time. The residual label is the Euclidean three-dimensional distance
between actual and reference positions at that time. It is not a closest-point
distance or an estimate of navigation error.

A reference marker requires aligned time within the reference interval; the
app does not extrapolate before or after the simulation. A trace or last point
can remain visible when new packets are absent, so always read sample age.
The vertical profile overlays reported altitude with the reference's
launch-relative up coordinate. Their comparison depends on using the correct
altitude convention.

\subsection{Altitude conventions}
For LIVE geographic-to-local conversion, configure an origin and a valid 3D
rocket GPS fix. Choose one verified interpretation:
\begin{longtable}{@{}L{.27\linewidth}L{.66\linewidth}@{}}
\toprule Mission choice & Conversion used for the rocket's absolute height \\\midrule\endhead
Unknown & Geographic trajectory position remains unavailable. Telemetry values still display. \\
GPS absolute ellipsoid & Use the reported GPS altitude directly. \\
GPS zeroed at launch origin & Add mission ellipsoid altitude to the reported GPS-relative altitude. \\
Filtered barometer relative to launch & Add mission ellipsoid altitude to the reported filtered barometric altitude. \\
\bottomrule
\end{longtable}
Mean-sea-level elevation is a separate simulation/weather input; it is not
substituted for ellipsoid height in the geographic transform. The legacy antenna
tracking convention is separate again (Section~\ref{sec:antenna}).

\subsection{Clock meanings}
There are three primary clocks. Device time is flight-computer uptime in
seconds after rollover extension. UTC is the packet's host reception time.
Monotonic host time measures freshness and session elapsed time. Flight time is
\[
 t_{\mathrm{flight}}=t_{\mathrm{device}}-t_0.
\]
In LIVE, observing a received Preflight-to-Flight transition establishes
\(t_0\). Connecting mid-flight does not establish liftoff time. A detected device
reboot clears time alignment and displayed history and stops tracking. The
Sessions action \ui{Set displayed sample to flight \(t=0\)} provides explicit
manual alignment.

The attitude display uses the legacy integrated rotations. In GNC, axes are
labeled X/roll, Y/pitch and Z/yaw. The attitude graphic and traces describe
the available legacy data; they do not imply a quaternion navigation estimator.

\section{Digital and Analog video}
Flight overview has independent \ui{Digital} and \ui{Analog} panels. Both receive
video from USB cameras; the labels describe the rocket-to-ground radio links.
In each panel choose \ui{Find}, select a listed USB camera/receiver and press
\ui{Start}. Each panel has its own \ui{Stop} and \ui{File\ldots} controls.
A selected camera is removed from the other dropdown; active and closing
captures also reserve their camera. Select the empty choice to release an
unused selection. Stopping or changing one stream leaves the other running. The app uses AVFoundation on Mac, DirectShow on Windows and
the existing Video4Linux branch on Linux. A local video file can also be selected.
Network-stream URLs are not a dedicated operator input in this version.

The displayed video is fitted into a \(960\times540\) RGB frame with its aspect
ratio preserved. Each stream has one FFmpeg input feeding its display and recording. Camera and
demo sources are compressed for recording; file inputs use stream-copy for the
recorded output. Audio is not captured. Files are stored as approximately
30-second Matroska segments with a CSV time index.

Each video status reports frame age and errors; after two seconds without a
recent frame it adds a stale/frozen indication. Stop closes capture but retains
the final displayed image with an explicit stopped label. A retained image is
not proof of a live stream.

Starting or stopping logging reopens both configured video pipelines so their
recording destinations change. A camera reconnect during a recording uses a
new segment directory rather than overwriting earlier footage. Video lifecycle
work occurs off the GUI thread, but a brief stream interruption is possible.

Replay uses host session timing plus the mission's video offset. A positive
offset delays video relative to telemetry; a negative offset advances it.
Each stream follows its own recorded starts, segments and gaps. A paused seek
decodes one frame per stream. Flight files include both feeds; old single-video
sessions appear in Digital. A panel's Stop pauses that replay feed until a seek
or playback restart. Compact displays can scroll to the lower flight plots.
Camera exposure timestamps are not
synchronized with onboard telemetry in this implementation.

\section{Antenna assembly and pointing}\label{sec:antenna}
\subsection{Viewing the physical model}
The drawing follows the supplied four-leg stand, azimuth turntable, elevation
supports and common antenna beam. The grid reflector, long Yagi and enclosed
Avenger XR18 share the elevation carrier; mounting centers are attached to the
beam concentric with the support pivots. The dish feed arm uses the dish's grid
material and lies on the dashed boresight. The Avenger enclosure represents the
black unit identified in the supplied photograph; no product photograph is
embedded in the renderer.

Drag to orbit the viewing camera, scroll to zoom, and double-click to reset the
perspective. Perspective, Front, Side and Top presets are available. Viewing
works with serial controls locked and never sends a command. All bearings use
the corrected chirality: \(0^\circ\) north, \(90^\circ\) east, \(180^\circ\) south,
\(270^\circ\) west.

\begin{figure}[htbp]\centering
\begin{tikzpicture}
\node[box,text width=35mm] (stand) at (0,0) {Four-leg stand\\fixed world frame};
\node[box,text width=35mm] (turn) at (5,0) {Turntable and supports\\azimuth rotation};
\node[accentbox,text width=35mm] (beam) at (10,0) {Common pivot beam\\azimuth and elevation};
\node[smallbox] (dish) at (5.7,-2) {Grid dish\\and aligned feed};
\node[smallbox] (yagi) at (10,-2) {Long Yagi\\beam attachment};
\node[smallbox] (helix) at (10,-3.8) {Avenger XR18\\beam attachment};
\draw[flow] (stand)--(turn);
\draw[flow] (turn)--(beam);
\draw[flow] (beam.south west)--(dish.north);
\draw[flow] (beam)--(yagi);
\draw[flow] (beam.east) -- ++(.6,0) |- (helix.east);
\node[captionnode,text width=42mm] at (.3,-2.8) {Orbit and zoom change the camera only. They do not change commanded azimuth or elevation.};
\end{tikzpicture}
\caption{Transform hierarchy of the procedural antenna model. Both outer antennas and the reflector belong to the same moving carrier; this diagram describes attachments, not dimensions or a photographic viewing direction.}
\label{fig:mount}
\end{figure}


\subsection{Manual controls and the meaning of ZERO}
Enter finite azimuth/elevation values and press \ui{Send}. Editing a text field
alone does not move the drawing or transmit. The graphic updates on successful
serial dispatch. The legacy manual path does not add mission calibration,
software envelope or cable-wrap prerequisites; stored mission envelope fields
are not enforcement for these buttons.

UP, DOWN, LEFT and RIGHT send the original firmware direction opcodes in
5-degree increments. These are direction packets, not host-generated absolute
targets. Before a known target exists, the app cannot show a reliable absolute
pose after a jog.

\ui{ZERO} sends the firmware's count/target reset at the current physical
position. It does not seek a home switch or mechanically return the mount to
north/horizontal. The drawing's last-sent label is a command estimate; measured
angles and board acknowledgment are unavailable.

\subsection{Ground-GPS freeze and automatic tracking}
With both boards connected and telemetry polling, press \ui{Send to AntPtr}.
The receiver latitude and longitude are frozen at the original five-decimal
rounding; its displayed altitude/fix are retained with that snapshot. The legacy
calculation uses ground altitude zero and rocket filtered barometric altitude.
Mission launch origin, ellipsoid/MSL settings and stored pointer offsets do not
replace this convention.

\ui{Track live rocket} computes targets at up to 5 Hz. It requires fresh LIVE
rocket data, usable rocket coordinates and a frozen ground position. The
geographic pointing helper rejects targets within 10 m and nearly overhead
targets with under 1 m horizontal separation. Stale data, reboot or connection
loss stops host tracking; resume explicitly after correcting the condition.
\ui{Unfreeze GPS} releases the stored origin and stops tracking.

\statusnote{Stop tracking}{This stops new targets and removes an unsent queued
target. A movement already commanded can continue. The old firmware has no
motion-stop opcode, watchdog, measured-pose report or mechanical-homing contract.}

\subsection{Virtual pointer connection and rehearsal}
The pointer selector also offers \ui{Virtual antenna pointer} in LIVE.
Alternatively, click \ui{Connect virtual pointer} on the antenna page.
This creates an in-process simulator and needs no ground station or physical
serial port. The connection, pose and banner explicitly identify simulated
motion; rocket controls remain locked until the real ground station connects.

Configure the mount under \ui{Mission > Configure > Antenna pointer}: WGS84
latitude, longitude and ellipsoid altitude, plus \ui{Antenna location established}.
Mission JSON and portable flight files retain the position. The physical-board
tracking path continues using the frozen receiver GPS from \ui{Send to AntPtr}.

Run OpenRocket or load a georeferenced, launch-relative ENU reference. Then use
\ui{Follow trajectory} on the antenna page. Pause/resume, Rewind, a time slider
and 0.25--4 times speed control an independent rehearsal clock. The gold target
and purple antenna marker in Flight overview identify the simulated geometry;
telemetry samples, flight counters and actual tracks are unchanged.

Send, the four direction controls, ZERO and their original keyboard shortcuts
actuate the simulated mount. Manual movement pauses playback. The displayed
pose approaches its target at illustrative limits of 90 degrees/s in azimuth
and 60 degrees/s in elevation. \ui{Hold virtual pointer} freezes the simulated
motion and clock. These rates are visualization assumptions, not measured mount
dynamics. Disconnect or a source-mode change closes the simulator. Editing the
antenna location or replacing the reference stops rehearsal until an explicit
restart. Connection and rehearsal state are not restored from a saved flight.
\begin{figure}[htbp]\centering
\begin{tikzpicture}
\node[box,text width=36mm] (reference) at (0,0) {OpenRocket reference\\ENU path + launch origin};
\node[accentbox,text width=36mm] (frame) at (5.15,0) {Antenna-local direction\\ENU via ECEF};
\node[box,text width=36mm] (pose) at (10.3,0) {Virtual pointer\\command + slew model};
\node[box,text width=36mm] (mission) at (0,-2) {Mission profile\\antenna lat/lon/height};
\node[box,text width=36mm] (clock) at (5.15,-2) {Independent clock\\play, pause, seek, speed};
\node[box,text width=36mm] (display) at (10.3,-2) {Simulated pose + markers\\telemetry unchanged};
\draw[flow] (reference) -- (frame);
\draw[flow] (frame) -- (pose);
\draw[flow] (mission.east) -- (frame.south west);
\draw[flow] (clock) -- (frame);
\draw[flow] (pose) -- (display);
\end{tikzpicture}
\caption{Virtual rehearsal uses the saved antenna position and its own clock. No physical serial transport participates in this path.}
\end{figure}


\section{Rocket commands, power and recovery}\label{sec:rocket}
\subsection{Telemetry subview}
The legacy Telemetry subview retains the original values: state, RSSIs, packet
age/number, filtered and maximum altitude, velocity, integrated rotations,
angle from vertical, temperature, GPS fix/position/heights/precision/satellites,
four servo drive/degree readouts and three cell voltages. Cells above 3.7 V use
the normal accent, those above 3.5 V the intermediate color, and the remainder
the low-value color. These cell cues are separate from the configurable
total-battery alert.

\subsection{Rocket commands subview}
\begin{longtable}{@{}L{.28\linewidth}L{.65\linewidth}@{}}
\toprule Control & Implemented behavior \\\midrule\endhead
Advance State & Requests current decoded state code plus one; retains the original confirmation dialog. The UI does not implement an independent flight-state machine. \\
Zero P/Y/Roll; Zero Alt; Zero Velo; Zero Servos & Sends the corresponding original opcode. These do not rewrite recorded historical samples. \\
PD Activate & Sends the legacy activation command. No new controller tuning interface is implemented. \\
Roll control angle; Airbrakes angle & Sends the entered finite angle as the original float32 payload. \\
VTX 1, 3, 5 or 8 W & Sends the legacy power-choice index. A selected button represents the request, not a radio acknowledgment. \\
\bottomrule
\end{longtable}
Ground commands require the ground board in LIVE even when polling is paused.
They are simulated in DEMO and disabled in REPLAY. The app does not require
fresh radio evidence before every uplink; the separate reception indicator must
be interpreted accordingly.

\subsection{Power subview}
Six converter requests and voltage/current readbacks are shown for 3, 3.3, 5,
7.4, 8.4 and 28 V. Changing an editable request sends the full six-bit mask.
The 3.3 V request stays on and is not editable, matching the old UI.
Total current and temperature are readouts, not switchable rails.
Readback state is derived from voltage proximity to the nominal rail, not from
the state of its checkbox.

The BMS table displays Enabled and Triggered rows for SCD, OCD2, OCD1, OCC,
COV, CUV and the two reserved bits. It is read-only in the original operator
workflow. A BMS command exists at protocol level but is not an additional
operator toggle.

\subsection{Recovery and pyro subview}
Six channels, numbered 0--5, have selections, continuity, armed/fired state and
resistance readouts. Select at least one channel before using ARM or FIRE.
Both retain the original confirmation behavior and encode the chosen channel
mask. The emergency actions PISTON, BP WELLS, TNDR DSNDR and DEPLOY ALL each
retain their own confirmation.

There is no added Disarm button in this subview, although the packet encoder
supports the legacy disarm opcode. Continuity and armed/fired displays come from
received packets; sending a request alone does not update them as confirmed
hardware state. Consult the team's operational procedures for physical recovery
system work; this report documents the existing UI rather than a firing sequence.

\section{Mission configuration and launch sites}\label{sec:mission}
\ui{Mission $\rightarrow$ Configure} opens a scrollable configuration dialog.
Saving validates the entries and applies them; Cancel leaves the current mission
unchanged. Stop logging before editing or loading another mission so a recording
is not redefined in progress. Fields include the mission name, origin-established
flag, location, two altitude datums, legacy height interpretation, canard count,
freshness/battery thresholds, launch rail, saved simulation index, random seed and
video offset. Appendix~\ref{app:mission} lists every stored field, including
legacy fields that are not active UI controls.

\subsection{Latitude/longitude, MGRS and Plus Codes}
All coordinate conversions are offline:
\begin{itemize}
\item \ui{Latitude / longitude}: decimal degrees, latitude \(-90\) to \(90\),
longitude \(-180\) to \(180\).
\item \ui{MGRS}: compact or spaced UTM/UPS grid reference. The preview identifies
the supplied grid resolution. Conversion uses the grid reference point, not a
newly estimated survey position.
\item \ui{Plus Code}: a full code resolves to its cell center. A short code
requires explicit nearby reference latitude and longitude; locality names are
not geocoded.
\end{itemize}
The preview shows the resolved canonical WGS84 coordinates before Save.
Changing only the display format preserves canonical coordinates and avoids
round-trip drift. A saved mission keeps the chosen format/code as entry
provenance alongside latitude and longitude.

\subsection{URRG preset}
Select \ui{URRG} from the launch-site preset menu. It sets the exact supplied
MGRS string \code{18TUN2061530290}, switches the entry format to MGRS and marks the
origin established. Its resolved coordinates are approximately
\(42.7041678^\circ\) latitude and \(-77.1902950^\circ\) longitude.

The preset supplies coordinates only. It does not infer launch elevation,
antenna position, wind or rail settings. Existing elevation values remain for
the operator to review. Manual coordinate edits clear the preset name to Custom;
changing the display format alone keeps URRG. Mission JSON and flight archives
preserve the site selection.

\subsection{Mission JSON versus portable flight files}
\ui{Save mission} writes configuration JSON. \ui{Open mission} restores it,
clears history/track and deselects the old reference. The JSON contains file
paths for a selected rocket and custom motors; it does not embed those files.
Loading mission JSON is not the same operation as opening a flight archive and
does not itself perform the archive workflow's full source disconnection.

Use \ui{Save flight} for a portable package containing the referenced assets
and available data. Keep the distinction in mind when moving work to another
machine.

\section{Wind and nominal OpenRocket simulation}\label{sec:simulation}
\subsection{Manual and imported wind}
The wind table uses height above launch in metres, speed in metres per second
and meteorological direction \emph{from} true north in degrees. Add or remove
layers and press \ui{Apply wind}. Editing a table cell alone does not update the
active mission profile.

Heights must be nonnegative and strictly increasing; there must be 1--200 layers.
Speeds must not exceed 150 m/s. A zero-speed direction has no physical effect.
Internally each row becomes east and north velocity components:
\[
 v_E=-s\sin\theta,\qquad v_N=-s\cos\theta .
\]
\ui{Import wind} and \ui{Export wind} use JSON with \code{layers} and a
\code{source} description. The included example shows the exact format.
Outside the supplied altitude coverage, the nearest wind layer is held.

\subsection{Online weather}
Enter an ISO 8601 time with a timezone, such as
\code{2026-09-21T16:00:00Z}, and the launch elevation above mean sea level.
\ui{Retrieve wind} requests an Open-Meteo pressure-level profile at the mission
latitude/longitude. The current implementation selects forecast or historical
forecast based on age, rejects requests more than 15 days ahead, and requires
an available provider hour within one hour of the selected time. These are app
rules, not a guarantee that the provider has every date/location.

Pressure levels are 1000, 925, 850, 700, 500, 300, 200 and 100 hPa. Missing levels
and below-launch levels are skipped; remaining layers are sorted.
Geopotential height is converted to geometric height, then the launch MSL
elevation is subtracted. The provider response, request, endpoint, retrieval time
and attribution remain in the mission for offline reuse. Review the resulting
coverage. If the launch coordinates change during retrieval, the UI rejects
applying that result to the new site.

\subsection{Running a reference}
\begin{enumerate}
\item Configure the launch origin and altitude conventions.
\item Choose the team's \code{.ork} file. It must contain a saved simulation with
a valid active motor configuration.
\item Select the saved simulation index in mission settings; index 0 means the
first saved simulation.
\item The Zephyrus MITRT N8406 curve is bundled automatically. For other custom
motors, use \ui{Select custom motor files} to provide matching \code{.eng} or
\code{.rse} curves. Cancel keeps the existing selection.
\item Review rail length, tilt from vertical, true-north heading, wind and seed.
Press \ui{Run nominal simulation}. \ui{Cancel simulation} terminates the worker.
\end{enumerate}
A fresh simulation uses the saved vehicle/motor configuration but explicit
app-supplied site, rail and wind settings. The bridge disables saved extensions
and the fork's global inertia override, uses the ISA atmosphere, zero wind
turbulence standard deviation and a fixed seed, and returns the first output
branch. The job also records branch count and warnings.

A missing thrust curve is named before simulation; an unpowered configuration
is rejected. The supplied Zephyrus test model previously failed because its
custom N8406 curve was absent. The bundled curve now matches its required motor
digest. The app copies model/motor inputs into each job and records hashes so
the run can be traced independently of original input paths.

\subsection{Importing a reference and interpreting results}
\ui{Import reference} on Flight overview accepts a CSV and a same-basename JSON
manifest. The CSV columns are \code{time_s,east_m,north_m,up_m}.
The manifest declares version 1, ENU, \code{m,s}, and
\code{origin:[latitude,longitude,ellipsoid_altitude]}; up is launch-relative.
A LIVE comparison requires a configured mission origin matching that reference
within \(10^{-6}\) degrees and 0.1 m altitude. Synthetic references are rejected
for LIVE comparison. Stop logging before replacing the recording's reference.

The saved trajectory is not a flight-control command, and the app does not
run an online controlled-vehicle simulation driven by each received packet.
For GNC assessment, compare the nominal baseline with actual data while retaining
the solver and unavailable-feedback limitations.

\section{Recorded Zephyrus demonstration}
Choose DEMO explicitly. GS2 is the default; GS1 and GS3 are alternative receivers
of the same launch, not consecutive segments. They are never concatenated or
deduplicated. Each plays according to its own receive-time deltas, keeping gaps,
original UTC, packet numbers and device time.

\Needspace{14\baselineskip}
\begin{longtable}{@{}L{.16\linewidth}L{.16\linewidth}L{.18\linewidth}L{.41\linewidth}@{}}
\toprule Receiver & Samples & Duration & Use and limitations \\\midrule\endhead
GS1 & 4,366 & 274.056 s & Original receiver record, with GPS outliers. \\
GS2 & 5,464 & 330.234 s & Default record; most samples and fewer major GPS corruptions. \\
GS3 & 4,939 & 332.725 s & Contains invalid ground GPS and anomalous rocket GPS/state readings; automatic simulated tracking is unavailable. \\
\bottomrule
\end{longtable}
The initial cue is five receive-time seconds before the first FLIGHT packet.
Use Play/Pause, the timeline, 0.25/0.5/1/2/4 times speed, \ui{Launch -5s} or
\ui{Full start}. Stop logging before seeking or changing the receiver log.
Pause and speed changes are allowed during a rehearsal recording.
At the end the final sample is retained; restarting returns to the launch cue.

The demo track uses horizontal GPS offsets from the first FLIGHT position plus
the reported barometric height. It is a local observed track, not a surveyed
ellipsoid/MSL trajectory. Missing 3D fixes and positions over 20 km horizontally
from its origin are omitted from the plot but kept in raw decoded values.
No interpolated packets, touchdown, actuator feedback or corrected state
sequence are invented. The records end in Main descent.

Both videos are labeled \ui{TEST PATTERN} (Digital motion pattern and Analog bars); no actual test-launch video was supplied.
The original 43 CSV fields remain in each decoded sample's provenance. Raw radio
frames were not supplied with these CSVs, so their radio checksums cannot be
revalidated. The app verifies the bundled CSV byte hashes on loading.

\section{Logging, sessions and portable flight files}\label{sec:persistence}
\subsection{Starting and stopping logging}
Start Logging creates a uniquely named session under the application data
directory's \code{sessions} folder. It records accepted samples, events,
raw serial bytes, both legacy CSVs and any recorded video. Logging runs
independently of UI redraws. A bounded recorder queue reports dropped records
and a fault if it overflows; loss is not silently labeled complete.

Stop Logging detaches the active recorder and finishes outstanding video and
file work in the background. The event timeline reports completion or error.
Closing the app also finalizes its workers. A crash can lose uncommitted data and
the active video segment; committed session data may still be replayable.

\subsection{Session-folder replay}
Use \ui{Open recorded session} to select a session folder. Playback begins paused.
Controls provide seek, a 0.1 s step, 0.25/0.5/1/2/4 times speed and manual time
alignment. Seeking reconstructs earlier events and up to 120 s of sample
history instead of retaining future pointer/time state. The complete session
remains indexed on disk even though visible plots are bounded.

\ui{Export samples to CSV} writes the richer normalized export, including elapsed
time and the complete sample JSON. This is distinct from the two original-format
CSV files written automatically by the recorder.

\ui{Import Zephyrus CSV} accepts the audited legacy layout, creates a converted
session in a chosen parent directory, retains the original source file and
opens it for replay. It preserves available extra columns; unsupported layouts
are rejected. Importing decoded CSV cannot recreate raw packet quality or radio
bytes.

\subsection{Save flight and Open flight}
Use \ui{File $\rightarrow$ Save flight} (Command-S on Mac) and
\ui{File $\rightarrow$ Open flight} (Command-O). Both are also available in Mission
and Sessions. A \code{.rktflight} archive contains the mission, embedded selected
rocket/custom-motor files, current normalized simulation reference, and available
session data. It is readable after the original model/motor paths disappear.

\begin{longtable}{@{}L{.28\linewidth}L{.65\linewidth}@{}}
\toprule State when saving & What is included \\\midrule\endhead
Logging active & Consistent committed telemetry/raw/CSV snapshot and finalized video segments. Logging continues. The current unfinished video segment is omitted and the session remains marked as a recording snapshot. \\
Logging just stopped & Last recording in the current source mode; save waits for pending recorder and video finalization. \\
REPLAY & Opened session and displayed replay position. \\
DEMO without a selected recorded session & Complete chosen GS dataset and current demo position; no invented camera footage or raw radio bytes. \\
LIVE without a recorded session & Only the displayed telemetry buffer, up to 6,000 samples, explicitly identified as a buffer export. \\
No available telemetry & Planning flight: mission, model/motors and any current simulation reference. \\
\bottomrule
\end{longtable}
Save again after stopping logging to include the final samples and video segment.
A flight archive carries normalized simulation output and metadata, but is not an
automatic copy of every file in the original Java job directory; keep that
directory for its complete worker log and inputs.

Opening requires logging to be stopped and a running simulation to finish or
be cancelled. The archive is validated before switching the current flight.
A session-bearing flight opens in paused REPLAY with both transports closed.
A planning flight opens in disconnected LIVE. Saved pointer calibration is
cleared; no hardware automatically reconnects.

Saving writes a temporary archive and atomically replaces the destination after
success; failure leaves an existing destination intact. Loading validates paths,
sizes, hashes and supported schemas. A damaged file reports an error instead of
silently supplying partial data. Extracted assets are retained under the app's
\code{flights} cache. The file format is versioned but not encrypted or
cryptographically authenticated.

\subsection{Storage and recovery}
The default root is Qt's per-user \code{AppLocalDataLocation}, under organization
\code{MITRocketTeam} and application \code{RocketGNCMonitor}. It normally lies
under macOS Application Support or Windows Local AppData. Use the paths displayed
by the app rather than assuming a fixed home directory. Developers can override
it with \code{--data-dir}.

A session contains \code{manifest.json}, \code{session.sqlite},
\code{raw.bin}, \code{telemetry.csv}, \code{telemetry_badpackets.csv}, and optional
reference, imported-source CSV and video directories.
Do not delete a crashed session's SQLite \code{-wal} file before recovery:
committed records can reside there. Incomplete sessions remain labeled
incomplete; a usable database is not a guarantee of complete footage or raw bytes.

\section{Alerts, diagnostics and troubleshooting}
Diagnostics shows packet counts, the current decoded sample, event messages and
recording information. Valid/rejected frames, discarded bytes and inferred
sequence gaps describe different conditions. Rejected payloads are isolated
from live sample displays and written to the bad-packet CSV when logging.

The LIVE alert manager tracks telemetry stale/absent, low total battery,
recording fault/loss and independently identified Digital/Analog video ended/stale. Battery must remain below threshold
for 1 s; clearing uses 0.3 V hysteresis. Video alerting has a 1 s dwell after an
error or a frame age over 2 s. Acknowledgment records an event and marks the alert;
it does not clear the underlying condition. The demo can show low original
battery values in readouts without those being a new live-hardware alarm.

\begin{longtable}{@{}L{.30\linewidth}L{.63\linewidth}@{}}
\toprule Symptom & Checks or explanation \\\midrule\endhead
Ground connected, no data & Start Polling. Check that the selected port is the telemetry board and that the rocket/receiver supplies valid frames. USB connection is not radio evidence. \\
Controls remain locked & Confirm mode LIVE and the relevant board connection. Pointer manual controls need the pointer board; rocket commands/logging need the ground board. \\
Camera absent or no image & Find cameras again, confirm the receiver enumerates as a camera, grant permission and inspect the video error. Another process may own the device. \\
Image visible but no motion & Read frame age and stopped/stale status; the last frame can be retained. DEMO uses a test pattern. \\
No actual trajectory & Establish launch origin, select the verified height convention and obtain valid 3D GPS. DEMO filters plot-only GPS outliers. \\
No same-time reference marker & Establish flight time alignment and ensure the displayed time falls within the reference interval. \\
Missing motor / unusable simulation & Select a saved powered configuration and matching custom curve. The N8406 is automatic. Review the named error and job's \code{worker.log}. \\
Tracking unavailable or stopped & Check both links, polling, frozen ground GPS, target freshness and valid coordinates. Restart deliberately after fixing the problem. \\
Graphic disagrees with mount & It is a last-sent command model, not feedback. ZERO resets counts at the physical position. Review firmware sign/mechanical reference separately. \\
Flight save has less video & Active snapshots omit the unfinished segment. Stop logging, wait for completion and save again. \\
Only a short live flight was saved & Logging was not active; the fallback saves only the bounded display buffer. Start Logging before acquisition for a full capture. \\
Flight file rejected & Inspect the error for version, checksum, missing asset, path, database or space problems. Keep the original and the session source; do not bypass validation. \\
Recording incomplete/loss & Inspect recorder error, dropped count and disk availability. Preserve SQLite/WAL and finalized video; neither raw nor CSV can recreate data never received or committed. \\
\bottomrule
\end{longtable}

\section{Keyboard shortcuts and operational limits}\label{sec:shortcuts}
Qt preserves the original shortcut strings. On Mac, Ctrl in those strings maps
to Command, and Alt maps to Option. The following actions work across workspaces;
their enabled state still follows the appropriate connection/mode.

\begin{longtable}{@{}L{.39\linewidth}L{.27\linewidth}L{.27\linewidth}@{}}
\toprule Action & Mac & Windows / Qt notation \\\midrule\endhead
Settings & Command-comma & Ctrl+comma \\
Open flight & Command-O & Ctrl+O \\
Save flight & Command-S & Ctrl+S \\
Refresh serial ports & Command-R & Ctrl+R \\
Ground connection toggle & Command-Option-C & Ctrl+Alt+C \\
Pointer connection toggle & Shift-Command-Option-C & Shift+Ctrl+Alt+C \\
Polling toggle & Command-Return & Ctrl+Return \\
Logging toggle & Command-L & Ctrl+L \\
Pointer UP / DOWN & Command-Up / Down & Ctrl+Up / Down \\
Pointer RIGHT / LEFT & Command-Right / Left & Ctrl+Right / Left \\
Pointer ZERO & Command-0 & Ctrl+0 \\
\bottomrule
\end{longtable}

Field acceptance still requires the actual boards, USB video receiver and mount,
representative data rates, reconnect/power-loss behavior and a launch-duration
soak. Three-axis control feedback, a qualified controlled-vehicle model, measured
pointer pose and firmware-level stop/acknowledgment are outside the presently
verified contract. The current development Mac distribution is ad-hoc signed;
production notarization and clean-machine acceptance remain separate release work.

\clearpage
\part{Codebase and developer documentation}
\section{Repository, entry points and source ownership}\label{sec:repo}
The application is self-contained under \code{ui_work/rocket_UI_v0}. Its
\code{src/rocket_gnc_monitor} package is the Python application;
\code{simulation_bridge/src/RocketBridge.java} is the Java integration layer.
The old PyQt5 ground-station folder is a behavioral reference, not a runtime
dependency. Production builds use bundled resources and do not need the original
telemetry, OpenRocket source or font-system paths.

The console entry point \code{rocket-gnc-monitor} and
\code{python -m rocket_gnc_monitor} both call \code{__main__.main}.
The packaged launcher calls that same function. It parses arguments, registers
Lucida Grande, names the Qt organization/application, resolves the per-user data
directory, constructs \code{MainWindow} and enters the Qt event loop.
The window owns one \code{Controller}; its close event invokes controller shutdown.

\begin{lstlisting}
rocket_UI_v0/
  src/rocket_gnc_monitor/    # application package
  simulation_bridge/src/    # Java bridge source
  resources/                # icon, fonts, demos, motors, examples
  docs/                     # in-app Markdown help and dependency notices
  tests/                    # unit, Qt, serial and engine integration tests
  tests/fixtures/           # captured legacy contract and small .ork fixture
  scripts/                  # setup, build and package verification
  packaging/                # PyInstaller spec, launcher, entitlements
  vendor/                   # staged engine, bridge, JRE, build manifests
  build/                    # generated jobs, logs and validation artifacts
  dist/                     # distributable app/DMG or folder/ZIP
  Makefile                  # shared native build entry points
  Makefile.macos / Makefile.windows
  pyproject.toml / uv.lock / requirements-build.txt
\end{lstlisting}

\subsection{Module map}
Paths in this report link to the reviewed source relative to the existing
repository layout. Use the named symbols rather than fixed line numbers as
maintenance anchors.

\begin{longtable}{@{}L{.29\linewidth}L{.64\linewidth}@{}}
\toprule Module & Responsibility and main seams \\\midrule\endhead
\src{src/rocket_gnc_monitor/__main__.py} & CLI/startup; Qt data path and fonts; startup, demo, flight and engine smoke modes. \\
\src{src/rocket_gnc_monitor/controller.py} & State owner and orchestration: connection generations, polling, command dispatch, tracking, mode changes, recording, replay, file jobs, alerts and shutdown. \\
\src{src/rocket_gnc_monitor/ui.py} & \code{MainWindow}, \code{MissionDialog}, workspace construction, shortcuts, file dialogs, result presentation, enabled states and plots. \\
\src{src/rocket_gnc_monitor/rocket_panel.py} & Complete legacy rocket Telemetry/Commands/Power/Recovery subviews and retained confirmations. \\
\src{src/rocket_gnc_monitor/settings.py} & Immutable \code{AppSettings}, validation, atomic per-machine persistence, runtime paths and JAR checks. \\
\src{src/rocket_gnc_monitor/settings_ui.py} & Modeless \code{SettingsDialog}, engine/folder pickers, defaults, validation and save/cancel. \\
\src{src/rocket_gnc_monitor/widgets.py} & Themes, \code{AttitudeView}, \code{MountView}, local geometry, material colors and camera/pose transforms. \\
\src{src/rocket_gnc_monitor/fonts.py} & Private font registration and global Qt font choice. \\
\src{src/rocket_gnc_monitor/domain.py} & \code{Mission}, \code{Sample}, finite checks, geographic transforms, pointing and wind validation. Contains a synthetic test-data helper. \\
\src{src/rocket_gnc_monitor/location.py} & Qt-independent coordinate conversion and \code{Location}. \\
\src{src/rocket_gnc_monitor/location_ui.py} & \code{LaunchLocation} widget; format/preset state, preview and edit-versus-display distinction. \\
\src{src/rocket_gnc_monitor/devices.py} & Port identity inventory and bounded independent \code{SerialWorker} threads. \\
\src{src/rocket_gnc_monitor/protocol.py} & Incremental \code{ZephyrusDecoder}, pointer packet builder and legacy route-validation helper. \\
\src{src/rocket_gnc_monitor/zephyrus.py} & Rocket command opcodes/packet building and canonical 43-column legacy CSV adapter. \\
\src{src/rocket_gnc_monitor/legacy_sample.py} & Restricted AST-based CSV literal parsing and full row-to-sample conversion. \\
\src{src/rocket_gnc_monitor/legacy.py} & Validated legacy CSV-to-session conversion, provenance and incomplete-import handling. \\
\src{src/rocket_gnc_monitor/demo.py} & Bundled CSV integrity, launch cue, independent receiver clocks and plot-only ENU frame. \\
\src{src/rocket_gnc_monitor/recording.py} & Asynchronous \code{SessionRecorder}, indexed read-only \code{SessionReader}, raw capture CRC reader. \\
\src{src/rocket_gnc_monitor/flight.py} & Portable archive snapshot, atomic save, validated extraction and \code{LoadedFlight}. \\
\src{src/rocket_gnc_monitor/media.py} & FFmpeg location, camera discovery, process arguments, frame queue, recording and orderly stop. \\
\src{src/rocket_gnc_monitor/virtual_pointer.py} & Port-free legacy packet simulator, bounded slew model and independent georeferenced \code{VirtualFlight} playback. \\
\src{src/rocket_gnc_monitor/trajectory.py} & \code{Trajectory}, online wind normalization, motor resources and isolated \code{SimulationJob}. \\
\src{src/rocket_gnc_monitor/__init__.py} & Package version identifier. \\
\src{simulation_bridge/src/RocketBridge.java} & Fork-specific OpenRocket loading, motor resolution, nominal options, trajectory/error publication and private preferences. \\
\bottomrule
\end{longtable}

\subsection{Dependencies and resource resolution}
Python is the main application language because Qt, serial I/O, numerical plots
and filesystem orchestration fit one runtime. Java is isolated to the existing
team OpenRocket fork. The GUI uses PySide6 Essentials and PyQtGraph with NumPy;
serial transport uses pyserial, video uses imageio-ffmpeg's native FFmpeg, and
location conversion uses \code{mgrs} and \code{openlocationcode}.

The project permits Python 3.12 through 3.14 in metadata; the build guidance and
tested environments center on 3.12--3.13. Dependency ranges live in
\src{pyproject.toml}; \src{uv.lock} and \src{requirements-build.txt} pin build
environments. Do not treat the range bounds as a tested matrix.

Source execution resolves resources from the repository root. Frozen execution
uses PyInstaller's \code{sys._MEIPASS}. FFmpeg is located through its wrapper;
Java/engine/bridge are resolved under \code{vendor}. A wheel or editable Python
package alone is not the complete production distribution: package the resource
tree and native runtimes through the supplied build workflow.

\section{Runtime architecture and state ownership}\label{sec:runtime}
\begin{figure}[htbp]\centering
\begin{tikzpicture}
\node[box,text width=38mm] (serial) at (0,0) {Serial worker threads\\decode and raw capture};
\node[accentbox,text width=40mm] (controller) at (5.4,0) {Qt main thread\\Controller and timer};
\node[box,text width=35mm] (ui) at (10.6,0) {MainWindow\\eight workspaces};
\node[box,text width=38mm] (recorder) at (0,-2.1) {Recorder thread\\SQLite, raw bytes, CSV};
\node[box,text width=40mm] (pool) at (5.4,-2.1) {Background job pool\\files, wind, simulation};
\node[box,text width=35mm] (video) at (10.6,-2.1) {Video workers\\Digital and Analog};
\node[smallbox] (java) at (5.4,-4.1) {Java subprocess\\OpenRocket bridge};
\draw[flow] (serial)--node[above,font=\footnotesize]{samples}(controller);
\draw[flow] (serial)--node[left,font=\footnotesize]{capture}(recorder);
\draw[flow] (controller)--node[above,font=\footnotesize]{signals}(ui);
\draw[flow] (controller)--(pool);
\draw[flow] (controller.south east)--(video.north west);
\draw[flow] (pool)--(java);
\draw[optional] (pool.south west) to[bend left=28] node[below,font=\footnotesize]{snapshot / restore} (recorder.south east);
\draw[optional] (video.east) to[bend right=26] (ui.east);
\node[captionnode,text width=35mm] at (10.6,-4.1) {Display frames may be dropped to keep the UI current. Recorder loss is reported separately.};
\end{tikzpicture}
\caption{Runtime ownership. Hardware reads, file compression and engine execution do not belong in Qt paint or input callbacks. LIVE accepted samples enter the recorder before the display queue.}
\label{fig:runtime}
\end{figure}

\subsection{Concurrency model}
Qt widgets and controller state are updated on the main thread. Serial workers
own their ports and decoders. The recorder has its own thread and SQLite
connection. General background jobs use a four-worker executor; video open/close
operations use one single-worker executor per stream to serialize its FFmpeg
lifecycle changes independently. Each decoder thread drains its FFmpeg frames,
with a separate stderr drain.

The controller's 30 ms timer drains up to 2,000 queued messages per tick,
advances DEMO/REPLAY, dispatches tracking when due, takes the most recent frame from each
video stream, collects completed futures and updates alerts. Playback delta is capped
at 0.3 s per tick, so a long blocked GUI does not produce an arbitrarily large
playback jump. This desktop scheduler is not a hard real-time control loop.

\begin{longtable}{@{}L{.30\linewidth}L{.21\linewidth}L{.42\linewidth}@{}}
\toprule Buffer / cadence & Bound & Consequence \\\midrule\endhead
Controller messages & 16,000 messages & UI overflow increments \code{ui_drops}; stale generations are ignored. \\
Displayed history and track & 6,000 each & Bounded UI memory; not the full session archive. \\
Plot time-series tail & 1,200 samples & Altitude/rate/rotation plots use a shorter recent tail. \\
UI refresh & About 10 Hz & \code{refresh} throttles updates to 0.1 s. Some tables refresh every 0.5 s. \\
Telemetry command queue & 64 commands & Bounded; each command expires after 0.5 s. \\
Pointer command queue & 1 command & No silent overwrite; a pending dispatch prevents another ordinary target. \\
Tracking & At most 5 Hz & New target only if the preceding one is not pending. \\
Recorder queue & 16,000 entries & LIVE overflow records loss/fault. File import can wait rather than drop rows. \\
Recorder commit & About every 0.25 s & Raw/CSV are flushed and fsynced before committing matching indexes/state. \\
Video display queue & 2 frames per stream & Old display frames are dropped for freshness, independently of FFmpeg recording. \\
Replay history reconstruction & Previous 120 s & Earlier samples remain on disk and can be sought explicitly. \\
\bottomrule
\end{longtable}

In LIVE, \code{Controller.enqueue} sends an accepted sample to the recorder
before putting it into the UI queue. Raw callbacks also reach the recorder
directly. Thus a dropped UI message need not mean the same sample was lost from
disk. Conversely, recorder overflow is an independent data-loss condition.
Do not merge these counters or move capture entirely behind UI repainting.

\subsection{Connection generations and asynchronous results}
A new serial connection receives a generation token. Worker messages are used
only if the current role still owns that generation and the mode is LIVE.
This prevents a late error/sample from an old worker from altering a new link.
Each video stream uses a separate generation and lifecycle executor to prevent
an old open result from replacing a newer selection. Camera reservations cover
pending opens and closes as well as running capture.

General jobs report through the controller's future list and \code{task_done}
signal. Weather/reference/simulation jobs retain source-mode identity; completed
work from a different mode is not silently installed. Simulation results are
also compared against the relevant mission inputs. The UI performs additional
weather site checks and reference-origin validation.

Flight operations use \code{flight_busy}. Saving captures copies of mutable
mission/reference/history state before background I/O. Opening stores the
connection generation and mission snapshot; if source/settings change while
loading, the result is not applied over the newer context.

\subsection{Mode transitions}
\code{switch_mode} validates the target and verifies a requested demo resource
before changing the current source. It stops logging, closes replay/serial/video
resources, cancels host tracking, clears histories and command state, resets time
and sample counters, clears pointer calibration and installs the new mode.
DEMO saves the prior reference for restoration when leaving.

\begin{figure}[htbp]\centering
\begin{tikzpicture}
\node[accentbox] (live) at (0,0) {LIVE\\explicit connection\\to hardware};
\node[box] (demo) at (5.2,0) {DEMO\\recorded CSV playback};
\node[box] (replay) at (10.4,0) {REPLAY\\session opens paused};
\node[box,text width=126mm] (gate) at (5.2,-2.1) {Transition boundary: stop logging and tracking; close serial/video; invalidate old workers; clear visible history; install the selected source};
\draw[flow] (live)--node[above,font=\footnotesize]{select}(demo);
\draw[flow] (demo)--node[above,font=\footnotesize]{load}(replay);
\draw[flow] (live)--(gate.north west);
\draw[flow] (demo)--(gate);
\draw[flow] (replay)--(gate.north east);
\node[captionnode,text width=125mm] at (5.2,-3.45) {A planning-only flight returns to disconnected LIVE. A session-bearing flight restores paused REPLAY. Source-mode names are independent of the rocket's onboard flight state.};
\end{tikzpicture}
\caption{Mode isolation. Other mode selections use the same transition boundary; the upper arrows illustrate common operator paths rather than the only permitted transitions.}
\end{figure}


\code{open_replay} validates the session's mission/reference before replacing
the current source, then enters REPLAY, restores the manifest, seeks to zero and
stays paused. \code{apply_flight} restores an extracted flight's mission,
reference and saved position; a planning-only file forces a clean LIVE state.
Neither path restores a physical connection.

\subsection{Lifecycle and shutdown}
\code{Controller.shutdown} is idempotent. It stops the Qt timer, stops serial
workers, closes video and waits for both video executors, cancels the active
simulation, closes recorder/reader resources and waits for background work.
Pending recording-close work is allowed to finish. Keep this ownership order
when introducing a new worker so it cannot write to an already closed resource.

\section{Data contracts, units and geographic frames}\label{sec:data}
\subsection{Mission and Sample}
\code{Mission} is a dataclass of persisted configuration. \code{validate}
checks coordinate/numeric ranges, supported location/height choices, actuator
count, seed/index, motor-path list and wind profile. \code{save} writes schema
version 1 through temporary-file replacement. \code{load} accepts version 1,
filters to known fields, uses dataclass defaults for absent fields and always
clears \code{pointer_calibrated}. There is no general multi-version migration
framework.

\code{Sample} is a dataclass representing one accepted or imported sample.
It contains device/receive times, source, packet number, navigation/kinematic
values, phase, actuator dictionary and extensible details. Optional missing
measurements use \code{None}; they are not replaced by zero. The live decoder
converts nonfinite floating sensor values to unavailable values. JSON output
rejects NaN. Appendices~\ref{app:mission} and \ref{app:sample} give all fields.

\code{details} carries original legacy fields, provenance, boot index and
quality annotations. Extend it conservatively: legacy rendering, CSV export,
record/replay and archive validation all consume samples. New actuator entries
should use explicit \code{demand}, \code{measured} and/or \code{drive} keys with
declared units, rather than renaming the four unknown legacy channels.

\subsection{WGS84 and local ENU}
\code{ecef} uses WGS84 semi-major axis \(a=6378137\) m and eccentricity squared
\(e^2=6.69437999014\times10^{-3}\).
For latitude \(\phi\), longitude \(\lambda\), ellipsoid height \(h\):
\[
 N(\phi)=\frac{a}{\sqrt{1-e^2\sin^2\phi}},\quad
 \boldsymbol{x}=\begin{bmatrix}
 (N+h)\cos\phi\cos\lambda\\
 (N+h)\cos\phi\sin\lambda\\
 (N(1-e^2)+h)\sin\phi
 \end{bmatrix}.
\]
\code{to_enu} subtracts the origin ECEF vector and applies the standard local
east/north/up rotation at that origin. The plotting contract is
\((E,N,U)\), with \(E\times N=U\).
\code{pointing} converts the local displacement to
\[
 \mathrm{az}=\operatorname{atan2}(E,N)\bmod 360^\circ,\qquad
 \mathrm{el}=\operatorname{atan2}\bigl(U,\sqrt{E^2+N^2}\bigr).
\]
The helper rejects the near-origin and nearly overhead cases described in Part I.
The live legacy pointer caller deliberately chooses its own origin/height
convention; it does not call the mission trajectory conversion as a substitute.

\begin{figure}[htbp]\centering
\begin{tikzpicture}[>=Latex]
\begin{scope}[shift={(0,0)}]
\draw[->,thick,cernblue] (0,0)--(2.5,0) node[right]{East};
\draw[->,thick,cernblue] (0,0)--(0,2.4) node[above]{North};
\draw[->,thick,cernblue] (0,0)--(-1.25,-.85) node[left]{Up};
\draw[->,teal,thick] (0,1.6) arc[start angle=90,end angle=0,radius=1.6];
\node[teal,font=\footnotesize] at (1.4,1.6) {azimuth};
\node[captionnode] at (.5,-1.6) {Navigation frame\\$E\times N=U$};
\end{scope}
\begin{scope}[shift={(7,0)}]
\draw[->,thick,cernblue] (0,0)--(2.5,0) node[right]{$+X$: west};
\draw[->,thick,cernblue] (0,0)--(0,2.4) node[above]{$+Y$: up};
\draw[->,thick,cernblue] (0,0)--(-1.5,-.9) node[left]{$+Z$: north};
\node[captionnode,text width=55mm] at (.4,-1.6) {Mount renderer coordinates\\East is $-X$, not $+X$};
\end{scope}
\end{tikzpicture}
\caption{The navigation and renderer bases are different but both preserve the intended chirality. The diagram is a projected axis sketch; up on the page is not a universal physical direction.}
\label{fig:frames}
\end{figure}


\subsection{Time normalization and reboot behavior}
The decoder extends a 32-bit millisecond rollover when a backward jump exceeds
\(2^{31}\). A smaller backward jump starts another boot epoch, resets rollover
offset and clears sequence-gap history. Sequence gaps use modulo-65536
differences and count missing values only for forward deltas between 2 and 32767.
The controller uses boot-index changes to invalidate alignment and old history.

For recording, elapsed time is the host monotonic receive time minus recorder
start, rounded to microseconds. Replay is driven by this elapsed column,
not by pretending onboard uptime began at zero. Events restore alignment and
last sent pointer state. Imported CSV uses original receive timestamps relative
to its first row. DEMO uses its receiver's own timestamps; it does not merge
the three clocks.

\section{Serial and protocol implementation}\label{sec:protocol}
\subsection{Port ownership and transport rules}
\code{ports} returns device, description, serial number and an identity derived
from VID/PID/serial when possible. \code{Controller.connect} enforces LIVE,
nonempty selection and distinct board ports, then constructs the worker.
Connecting does not start telemetry polling or transmit a setup command.

The worker opens pyserial with the role-specific read timeout and a 0.2 s write
timeout. Queued writes call \code{cancel_read} where available to wake the worker.
Only the telemetry role has the polling gate. Pointer RX is logged as bytes and
hex status, not decoded as authoritative measured angles.
Each write checks expiration, records the actually written prefix, and emits
\code{sent} only if the complete payload was written. Errors close the port;
there is no retry loop.

\subsection{Telemetry framing and checksum}
A telemetry frame is 144 bytes:
\begin{lstlisting}
offset 0..1      AB AB synchronizer
offset 2..129    128-byte Zephyrus payload P
offset 130..143 14-byte ground receiver trailer G

P[127] = sum(P[0:127]) modulo 256
\end{lstlisting}
The incremental parser retains partial frames and a trailing potential sync
byte. After a bad checksum it advances one byte and rescans, so an overlapping
valid synchronizer is not thrown away. It accepts concatenated reads and
reports accepted, rejected, discarded-byte and inferred-gap counters.
The additive checksum covers the payload only, is weak error detection and
does not authenticate a sender.

Payload offsets are relative to \(P[0]\); multibyte fields below are little
endian unless stated otherwise.
\begin{longtable}{@{}L{.14\linewidth}L{.31\linewidth}L{.48\linewidth}@{}}
\toprule Bytes & Field & Decode / interpretation \\\midrule\endhead
0--1 & Pyro continuity & Six two-bit states from the packed 16-bit value. \\
2, 3 & Armed, fired & Channel bit masks; channels 0--5. \\
4--9 & Pyro resistance & Six unsigned bytes divided by 10. \\
10--15 & Servo drives & Four packed 12-bit values, retained as legacy drives. \\
16--24 & Acceleration XYZ & Three signed 24-bit integers divided by 12800 and multiplied by 9.80665; m/s$^2$. \\
25--30 & Gyro XYZ & Signed int16 values times 0.03051757812; legacy Y sign inverted; degrees/s. \\
31 & Rocket GPS fix & Raw fix code; 3D trajectory requires at least 3. \\
32--39 & Latitude, longitude & Two signed int32 values times $10^{-7}$ degrees. Invalid bounds clear coordinates/fix. \\
40--43 & GPS height & Float32 metres; interpretation is mission-dependent. \\
44--51 & Horizontal/vertical accuracy & Two uint32 values divided by 1000. \\
52 & Satellites & Unsigned byte. \\
53--58 & Raw pressure, temperature & Two 24-bit raw words; temperature calibration is in the legacy adapter. \\
59--62 & Filtered altitude & Float32 metres. \\
63 & State & 0 ground testing, 1 preflight, 2 flight, 3 post-apogee, 4 main, 5 end. \\
64--75 & Integrated rotations & Three float32 values at 64, 68, 72; rendered roll/pitch/yaw. \\
76--79 & Maximum heights & Barometer then GPS uint16. \\
80--83 & Onboard clock & uint32 milliseconds uptime. \\
84--85 & Packet sequence & uint16 modulo counter. \\
86 & Onboard RSSI & Signed byte minus 99. \\
87--92 & Cell voltages & Three signed int16 values divided by 1000. \\
93--94 & Total current & Signed int16 divided by $-1000$. \\
95--97 & BMS & Temperature byte divided by 2; protection status and enabled bits. \\
98--109 & Six rail voltages & Signed int16 times 0.0016. \\
110--121 & Six rail currents & Signed int16 times 0.000625. \\
122--125 & Velocity & Float32 metres/s. \\
126 & Unmapped byte & Not interpreted by the present decoder. \\
127 & Payload checksum & Sum modulo 256 of bytes 0--126. \\
\bottomrule
\end{longtable}
Trailer \(G\) is RSSI at 0 (signed byte minus 99), fix at 1, signed int32
latitude at 2, longitude at 6 (both \(10^{-7}\) degrees), and signed int32
height at 10. The legacy adapter divides receiver height by 1000 even though the
audited firmware supplies metres. Preserve this discrepancy explicitly; the
legacy pointer origin instead uses height zero.

\subsection{Pointer command packets}
\begin{lstlisting}
byte 0       AA
byte 1       opcode
bytes 2..5   azimuth float32 LE (opcode 0), otherwise zero
bytes 6..9   elevation float32 LE (opcode 0), otherwise zero
byte 10      sum(bytes 1..9) modulo 256
\end{lstlisting}
Opcodes are 0 absolute, 1 UP, 2 DOWN, 3 LEFT, 4 RIGHT and 5 ZERO.
The legacy firmware negates internal azimuth and uses a shortest-route policy;
the host preserves original direction opcodes rather than reinterpreting them.
Absolute finite values are serialized without invoking the old
\code{validate_route} helper.

\code{validate_route} and mission envelope/calibration fields remain in source
for compatibility and tests of that helper. They are not called by the current
manual control/tracking dispatch path. Do not document them as an active
mechanical interlock or silently enable them when modifying the legacy UI.

\subsection{Rocket uplink packets}
\begin{lstlisting}
byte 0       AA
bytes 1..12  zero/reserved, except:
             byte 12 integer argument, or
             bytes 9..12 float32 LE angle
byte 13      command opcode
bytes 14..15 sum(bytes 1..13) modulo 65536, BIG endian
\end{lstlisting}
The mixed endian checksum is deliberate and covered by captured vectors.
\begin{longtable}{@{}L{.15\linewidth}L{.40\linewidth}L{.38\linewidth}@{}}
\toprule Opcode & Encoder command & Argument \\\midrule\endhead
01 / 06 / 07 & Arm / fire / disarm pyros & Nonempty bitmask of channels 0--5. \\
02 & Advance state & Requested state code; UI uses current + 1. \\
03 / 05 & Airbrakes / roll servo & Finite float32 angle. \\
04 & PD activate & No parameter. \\
08 / 09 & Zero servos / P-Y-roll & No parameter; zero-roll API aliases 09. \\
0A / 0B & Zero altitude / velocity & No parameter. \\
10 / 11 / 12 / 13 & Piston / BP wells / tender descender / all & No parameter. \\
14 & VTX power & Index 0--3 for 1/3/5/8 W. \\
15 & Converter requests & Six ordered Boolean request bits. \\
16 & BMS protections & Boolean encoded 0/1; API exists, operator status is read-only. \\
\bottomrule
\end{longtable}
\code{rocket_packet} validates command/value types. UI confirmation and transport
authorization remain separate layers: construct bytes without side effects,
confirm the original sensitive actions in \code{RocketPanel}, then dispatch
through the controller.

\section{Recording, replay and portable archive internals}\label{sec:storage}
\subsection{Session writer and durability}
\code{SessionRecorder} creates
\code{session_YYYYMMDD_HHMMSS_<suffix>} and writes an initial manifest with
\code{complete:false}. It enqueues sample dictionaries, events and role-tagged
raw byte chunks. One writer thread owns all session output. SQLite uses WAL and
\code{synchronous=FULL}; ordinary commits occur about every 0.25 s and at closure.

\begin{lstlisting}[language=SQL]
CREATE TABLE samples(
  id INTEGER PRIMARY KEY, elapsed REAL, t REAL, data TEXT);
CREATE INDEX sample_time ON samples(elapsed);
CREATE TABLE events(
  id INTEGER PRIMARY KEY, elapsed REAL, data TEXT);
CREATE INDEX event_time ON events(elapsed);
CREATE TABLE raw_index(
  id INTEGER PRIMARY KEY, elapsed REAL, role TEXT,
  offset INTEGER, length INTEGER);
CREATE TABLE snapshot_state(
  id INTEGER PRIMARY KEY CHECK(id=1), data TEXT);
\end{lstlisting}
\code{samples.data} is JSON from \code{Sample.to_dict};
\code{events.data} is \code{{"name":..., "data":...}}.
A raw index offset points to the chunk header in \code{raw.bin}; its length is
the payload length, not the whole chunk length.

Before a database commit, raw bytes and both CSV outputs are flushed and fsynced.
The same SQLite transaction stores a \code{snapshot_state} row with matching
raw/CSV byte lengths, manifest snapshot, duration, error and dropped count.
A successful final close checkpoints WAL and updates the manifest with completion,
duration and error/loss information. Queue overflow marks loss; a close timeout
or writer exception leaves an incomplete/error condition.

The recorder's own decoder watches raw telemetry RX to produce the rejected
packet CSV. Bad packets do not become \code{samples} rows. Good and bad CSV field
order comes from \code{CSV_FIELDS}; the adapter retains the original padding and
calibrations while fixing the documented old bad-packet-only inconsistencies.

\subsection{Raw capture format}
Each raw chunk has a 25-byte little-endian header with
\code{struct.Struct("<4sddBI")}:
\begin{longtable}{@{}L{.17\linewidth}L{.23\linewidth}L{.53\linewidth}@{}}
\toprule Offset & Type & Meaning \\\midrule\endhead
0--3 & 4 bytes & Magic \code{RGM1}. \\
4--11 & float64 & Host session elapsed seconds. \\
12--19 & float64 & UTC Unix seconds. \\
20 & uint8 & Role: telemetry RX 1, pointer RX 2, pointer TX 3, telemetry/rocket uplink TX 4. \\
21--24 & uint32 & Payload byte count. \\
25 onward & Byte payload & Exactly the declared number of raw serial bytes. \\
Following payload & uint32 LE & CRC32 of header plus payload. \\
\bottomrule
\end{longtable}
\code{read_raw} yields complete CRC-checked chunks and rejects wrong magic,
payloads over 4 MiB, truncated headers/payloads and CRC mismatch. This capture CRC
protects stored chunks; it does not add authentication to the radio or repair
the unchecksummed receiver trailer.

\subsection{Replay reader}
\code{SessionReader} opens SQLite read-only, reads manifest version 1, gets sample
count and maximum sample elapsed time, and runs \code{PRAGMA quick_check}.
\code{at(t)} returns the latest sample at or before \(t\);
\code{between(a,b)} returns \(a<elapsed\le b\), ordered by elapsed and id,
with a default 20,000-row limit. \code{events(t,begin)} uses the same temporal
selection convention.

\code{seek} clamps to the session duration, clears future state, resets initial
time alignment, applies prior events and reconstructs a 120 s history window.
Received monotonic times are refreshed for display. Normal play advances the
cursor at the selected speed, applies events and samples in the traversed
interval, and pauses at the end. Events currently restore flight-zero and
sent-pointer state; do not assume every console string is executable replay logic.

\subsection{Flight archive layout}
The portable extension is \code{.rktflight}; its contents are ZIP64 with a
\code{flight.json} inventory. Files normally use DEFLATE level 3, while MKV is
stored without another compression layer.

\begin{lstlisting}
flight.json                 # format/version, source/scope, position, inventory
mission.json                # version 1; asset paths relative to this archive
assets/rocket.ork            # optional selected model
assets/motor_000.eng         # optional selected curves (.eng or .rse)
reference.csv
reference.json              # optional current normalized reference
session/manifest.json       # only if a session was exported
session/session.sqlite      # standalone staged database, no -wal/-shm
session/raw.bin
session/telemetry.csv
session/telemetry_badpackets.csv
session/legacy-source.csv   # if available from legacy import
session/reference.*         # if present
session/video*/segments.csv
session/video*/segment_*.mkv
\end{lstlisting}
The inventory records byte count and SHA-256 for every payload file; it does not
recursively checksum itself. Top-level keys include
\code{format:"rocket-gnc-flight"}, \code{schema_version:1},
app version, save UTC, source mode, export scope, sample count, optional session
path, playback position, flight-zero/alignment and video-scope description.

Selected user motor files are embedded under generated portable names. The
app's own standard/bundled motor library is part of the runtime distribution;
it is not duplicated in every flight. Original job hashes and request metadata
may remain in the reference manifest as provenance; the live model/motor paths
are rebound to extracted assets.

\subsection{Active snapshot consistency}
\begin{figure}[htbp]\centering
\begin{tikzpicture}
\node[box,text width=37mm] (write) at (0,0) {Recorder keeps appending\\raw bytes and CSV};
\node[box,text width=37mm] (commit) at (5.2,0) {Flush files, then commit\\rows and prefix lengths};
\node[accentbox,text width=37mm] (backup) at (10.4,0) {Online SQLite backup\\one committed snapshot};
\node[box,text width=37mm] (prefix) at (10.4,-2) {Copy matching prefixes\\and finalized video};
\node[box,text width=37mm] (zip) at (5.2,-2) {Hash inventory and ZIP\\standalone database};
\node[box,text width=37mm] (publish) at (0,-2) {Atomic destination replace\\existing file kept on failure};
\draw[flow] (write)--(commit);
\draw[flow] (commit)--(backup);
\draw[flow] (backup)--(prefix);
\draw[flow] (prefix)--(zip);
\draw[flow] (zip)--(publish);
\draw[optional] (commit.north) to[bend right=35] node[above,font=\footnotesize]{next acquisition batch} (write.north);
\end{tikzpicture}
\caption{Save during logging. The snapshot boundary is a committed database transaction with matching file lengths, not an uncontrolled copy of an active SQLite file.}
\label{fig:snapshot}
\end{figure}

\code{snapshot_session} uses SQLite's online backup API to capture a committed
database without pausing the serial writer. It then reads the copied
\code{snapshot_state} to find exactly which raw/CSV prefixes belong to that
database version. Later appends to the source do not change those prefixes.
An early save waits briefly for the first committed state. Older stable sessions
without snapshot state use their available files and retain their completion
status.

Video is copied from finalized entries in the segment CSV; a partial trailing
index line is ignored for active capture, and segments beyond the snapshot
duration are omitted. After logging stops, the controller waits for pending
video/recorder closure before saving. No unfinished camera file is claimed as
complete footage.

The staged database is checkpointed and changed to DELETE journal mode once all
staging readers close. Only this owned staging copy is changed; the active
source database stays in WAL mode. This avoids shipping platform-specific
shared-memory state. A temporary archive is fsynced and replaced atomically at
the requested destination.

\subsection{Archive validation and application}
\code{load_flight} extracts into a new UUID directory under the owned flight
cache. It rejects unsupported versions, encryption, linked/directory/duplicate
case-folded entries, absolute/traversal/backslash/colon/control-character paths,
mismatched inventories, declared sizes and content hashes. Loading limits are
64 GiB of expanded entries, 100,000 entries and 16 MiB per metadata object; free
space is checked before extraction.

Mission asset paths must resolve to included assets. Mission/reference data,
SQLite integrity, sample/event structure and video segment indexes are validated
before the controller applies a \code{LoadedFlight}. On failure, the new
extraction directory is removed. Parsing uses JSON and ordinary media/model
formats, not pickle or evaluation of archive-supplied Python.

The archive inventory is corruption detection, not a signature or trust identity.
A deliberately altered payload with a recomputed hash is still subject to
schema checks, but is not proven to have come from the team. File limits and
validation should remain aligned when evolving save and load behavior.

\Needspace{24\baselineskip}
\subsection{Read-only analysis example}
Run analysis with the app's development environment so its domain contracts
match the recording:
\begin{lstlisting}[language=Python]
from pathlib import Path
from rocket_gnc_monitor.recording import SessionReader, read_raw

session = Path("/absolute/path/to/session_folder")
reader = SessionReader(session)
try:
    print(reader.count, reader.duration, reader.manifest["complete"])
    for elapsed, sample in reader.between(-1, min(10, reader.duration)):
        print(elapsed, sample.t, sample.altitude)
    reader.export_csv(session.parent / "analysis-export.csv")
finally:
    reader.close()

for elapsed, utc, role, payload in read_raw(session / "raw.bin"):
    print(elapsed, role, len(payload))
\end{lstlisting}
Do not feed replay events to a serial worker. Reading archive/session data is
separate from authorizing hardware commands.

\section{Legacy import and demo provenance}\label{sec:legacy}
\code{numeric_list} uses a small AST visitor. It accepts numeric/Boolean literals,
bounded list/tuple forms, unary signs and known one-argument NumPy numeric
wrappers, plus known Boolean names. It rejects arbitrary calls and attributes;
it never invokes \code{eval}. Numeric-array text is limited to 4,096 characters
and arrays to 32 elements.

\code{import_legacy} rejects files over 200 MB and known bad-packet filenames,
requires the audited core headers, checks monotonic receive timestamps and
refuses empty datasets. File import can block on the recorder queue to preserve
all rows. It stores source bytes and hash, records observed imported flight
transitions, and explicitly marks raw packets unavailable. A failed import
retains an incomplete/error manifest rather than presenting a completed session.

\code{DemoFlight} verifies the uncompressed SHA-256 and row count from
\code{resources/demo/manifest.json}; gzip is lossless packaging. It computes
receiver-relative times, locates the first FLIGHT row and exposes a launch cue
and indexed samples. Each sample contains the original row/UTC, station,
source file/hash and local-frame note. Plot-only rejection of invalid/far GPS
does not erase the original values.

The functions \code{domain.demo_sample} and \code{Trajectory.demo} remain useful
synthetic fixtures and examples. They do not drive the current operator DEMO
telemetry. Avoid reconnecting these helpers to the source selector when adding
features; the user-visible default is the actual GS2 log.

\section{Video implementation and synchronization}\label{sec:media}
\code{camera_devices} invokes bundled FFmpeg to enumerate AVFoundation or
DirectShow devices (and Linux video devices in its existing branch).
\code{VideoWorker.arguments} builds an argument list, without a shell, for camera,
file or test-pattern input. It maps only video and disables audio.

The display pipe is \code{rgb24} at \(960\times540\). The decoder reads exactly
one full frame before queueing it; incomplete EOF frames are discarded.
The two-frame queue trades display history for low latency. Recording is a
separate FFmpeg output from the same input: MPEG-4 for live/test sources,
stream-copy for local files, Matroska segments and \code{segments.csv}.
The display scaling does not imply the stored camera stream has the same
dimensions.

\code{Controller.video_streams} maps \code{digital} and \code{analog} to
\code{VideoStream} state: worker, backend, configuration, generation, reserved
cameras, replay key/pause state, lifecycle future and open error. \code{video_frame} carries
the stream identifier and bytes; \code{video_reset} clears only that panel.
The legacy \code{video}/\code{video_config} properties and \code{frame} signal
remain first-stream aliases.

\code{start_video(..., stream="digital")} increments that stream's token and
schedules work on its single-worker executor. Duplicate camera sources are
rejected before altering either stream, including a camera whose prior handle
is still closing. \code{stop_video(stream)} closes one feed; omitting the stream
closes both. The UI filters camera selections and uses \code{serial_device_key}
to exclude active board ports and macOS aliases from the other serial dropdown. \code{stop} first writes \code{q} and continues
draining output so FFmpeg can finish the trailer/index, then uses bounded
terminate/kill fallback. Keep stderr draining independent so logs cannot block
frame acquisition.

A \code{Video recording started} event records session elapsed time, source kind,
video directory and \code{stream}. Every recording start allocates a unique
\code{video_<stream>_<uuid>} directory, even when opens are queued rapidly.
Stopping logging waits for both lifecycle futures before closing the recorder,
including a feed manually stopped just beforehand. Old events without a stream
identifier map to Digital. Replay finds the eligible recording segment at
\[
 t_{\mathrm{video}}=t_{\mathrm{replay}}-t_{\mathrm{start}}-\Delta_{\mathrm{offset}}.
\]
It chooses each stream's segment whose interval contains that time and seeks
relative to its start. A manually stopped replay feed stays stopped until an
explicit seek or playback restart. Playing checks segment transitions every 0.2 s; a replay
key avoids reopening an unchanged segment. The reader's media events and
segment indexes are cached so these checks do not repeatedly scan the event log. Gaps stop only the affected feed. Positive offset therefore delays video. Paused replay uses
\code{single_frame}; playback uses the requested speed. This is host timing
alignment, not camera exposure or hardware clock synchronization.

\section{Weather and OpenRocket bridge contracts}\label{sec:bridge}
\subsection{Weather normalization}
\code{weather_profile} parses timezone-aware ISO timestamps and normalizes to
UTC. It caps response bytes at 4,000,000 and uses a 15 s network timeout. The
nearest returned hour must lie within 3,600 s. It transforms geopotential height
\(H\) to geometric height using \(R=6356766\) m:
\[
 h_{\mathrm{AGL}}=\frac{RH}{R-H}-h_{\mathrm{launch,MSL}}.
\]
The returned object contains \code{layers}, \code{raw}, \code{source},
\code{request}, \code{endpoint} and \code{retrieved_utc}. Network work stays in the
job pool. The mission stores that object in \code{weather_raw}. Manual profiles
use exactly the same validated east/north layer representation.

\subsection{Trajectory contract}
A trajectory is an \(N\times4\) NumPy array \((t,E,N,U)\) plus JSON metadata.
It requires 2--200,000 rows, finite values and strictly increasing time.
Loading caps the CSV at 40 MB. The manifest must declare version 1, ENU and
\code{m,s}; a nonsynthetic reference needs a valid three-value geographic origin.
\code{at(t)} interpolates each spatial coordinate linearly within coverage and
returns unavailable outside it. Save produces CSV and a same-basename JSON.

\subsection{Simulation job boundary}
The Python job copies the model and all explicit/bundled custom curves, writes
\code{request.json}, launches the private Java executable and captures its output
in \code{worker.log}. The request contains schema version, copied asset paths,
bundled motor count, simulation index, launch latitude/longitude, MSL altitude,
ellipsoid origin altitude, rail settings, seed, wind layers, output directory
and nominal/inertia flags.

\begin{figure}[htbp]\centering
\begin{tikzpicture}
\node[box,text width=36mm] (inputs) at (0,0) {Mission + ORK + motors\\wind, rail and seed};
\node[box,text width=36mm] (job) at (5.15,0) {Python job directory\\copied inputs + request};
\node[accentbox,text width=36mm] (java) at (10.3,0) {Private Java process\\RocketBridge + team engine};
\node[box,text width=36mm] (result) at (10.3,-2) {Trajectory CSV + JSON\\or structured error};
\node[box,text width=36mm] (validate) at (5.15,-2) {Validate trajectory\\attach hashes and provenance};
\node[box,text width=36mm] (view) at (0,-2) {Qt result handling\\reference-origin checks};
\draw[flow] (inputs)--(job);
\draw[flow] (job)--(java);
\draw[flow] (java)--(result);
\draw[flow] (result)--(validate);
\draw[flow] (validate)--(view);
\node[captionnode,text width=130mm] at (5.15,-3.3) {Cancellation and worker limits do not stop telemetry acquisition. Simulation results remain nominal and are never translated directly into physical control commands.};
\end{tikzpicture}
\caption{The simulation boundary uses files and a subprocess, keeping the JVM outside the Qt process.}
\end{figure}


Limits are 40 MB for the model file, 100 MB expanded ORK archive contents,
10 MB per motor curve, a configurable worker time limit (120 s by default) and 20 MB of worker log.
Cancellation sets an event and terminates the child; timeout/oversized-log paths
also stop it. Result collection checks return status, prefers a structured
\code{error.json} message and otherwise includes a bounded log tail.

The bridge runs headless with private in-memory preferences. It loads the team's
motor database, adds hash-verified bundled curves and explicit custom curves,
loads the ORK document and verifies the selected saved configuration. A missing
motor warning becomes a named error before simulation. A configuration without
an active motor is rejected.

The fork-specific nominal setup disables saved simulation extensions and its
global inertia override. It sets fixed rail direction rather than
\emph{launch into wind}, a 0.05 s step setting, a 600 s simulation maximum,
ISA atmosphere, multilevel AGL wind and zero turbulence standard deviation.
The Java wind direction is \emph{toward},
\(\operatorname{atan2}(v_E,v_N)\), in contrast to the operator's meteorological
\emph{from} input.

The first branch contributes time, X/east, Y/north and altitude/up.
Nonfinite/nonincreasing rows are skipped; fewer than two rows is an error.
Output metadata records branch count, motors and their engine digests, saved
simulation name, warnings, nominal qualification flags and solver description.
Python adds engine/model/motor hashes and the exact request/job directory.
The bridge exits explicitly after success so engine loader threads cannot keep
the worker alive.

\subsection{Reproducibility and interpretation}
The bundled MITRT N8406 curve is \code{resources/motors/TestLaunch_v14.eng}.
Its engine digest is
\code{057a820844b93fe3307401623a92d4c6}.
Byte integrity and model-digest matching are separate checks: identical naming
alone is insufficient.

A fixed seed and copied inputs support reproducibility on the pinned engine.
Results do not automatically transfer to a different OpenRocket binary or
qualified controlled-flight model. The package retains the complete team engine;
a random upstream JAR is not a drop-in replacement for fork-specific fields.
The exact source revision corresponding to the supplied engine binary remains
a release provenance item in the existing build metadata.

\section{UI, font and antenna graphics maintenance}\label{sec:graphics}
\subsection{UI responsibilities}
Workspace builders in \code{MainWindow} construct widgets; event handlers call
the controller or submit background jobs; \code{refresh} presents snapshots.
\code{refresh_connections} and \code{RocketPanel.set_controls_enabled} express
the existing board/mode rules. \code{guard} reports exceptions in the event area,
status bar and warning dialog. Do not swallow a rejected operation by leaving
a success-looking UI state.

\code{install_shortcuts} creates each QAction once. The original ten serial
actions coexist with standard Open/Save flight actions. Binding signals again
during periodic refresh would duplicate transmissions and must be avoided.
The daylight toggle updates the shared palette, stylesheet and plot/widget
colors. New custom-drawn labels should use the same font and palette conventions.

The global application font is Lucida Grande, registered from the packaged
font collection by \code{configure_fonts}. The collection is loaded privately;
the app does not install a system font. Plot axis labels and painter text use
the same family. The report itself uses Latin Modern to match the requested
LaTeX report style; these are separate presentation choices.

\subsection{Where the antenna's physical description is defined}
The physical drawing is procedural code inside
\src{src/rocket_gnc_monitor/widgets.py}, specifically
\code{MountView.paintEvent}. There is no separate CAD/mesh file.
Local helper functions \code{box}, \code{rod} and \code{face} construct polygons;
\code{palette} defines timber, frame, metal, grid, dark enclosure and other
material colors. Faces are depth-sorted and projected with QPainter, so the
drawing does not require an OpenGL/WebGL viewer.

The local transform functions deliberately separate attachment levels:
\begin{itemize}
\item \code{identity}: fixed stand/legs.
\item \code{turn}: azimuth rotation for turntable and supporting structure.
\item \code{carrier}: azimuth times elevation rotation, then translation to
the pivot at renderer position \((0,3.28,0)\).
\item \code{avenger}: carrier-local offset \((-1.03,0,0.16)\) for the enclosed
Avenger assembly.
\end{itemize}
The common beam is centered at carrier-local zero with extents
\(2.65\times0.10\times0.10\) in illustrative model units.
The Yagi boom is at local \(x=1.03\), running approximately from \(z=-0.35\)
to \(3.55\), with 21 crossed element positions. The grid is sampled procedurally;
the feed extends along local \(+Z\) from \(z=0.27\) to \(1.04\), in grid material.
The dashed boresight continues on the same axis. These values define appearance,
not fabrication dimensions.

\code{mount_rotations} and \code{MOUNT_COMPASS} implement the corrected basis:
renderer \(+Y\) up, \(+Z\) north, \(-X\) east. Do not flip only the compass letters
to compensate for a reflection elsewhere. The regression checks
\(E\times N=U\), headings at the cardinal azimuths and determinant \(+1\).

Camera yaw/pitch start at \(72^\circ/20^\circ\). Dragging changes those at
0.45 degrees per pixel. Zoom is clamped to 0.65--2.5. The analytic right/up basis
stays defined through top/bottom views and full orbit. Camera state and
\code{set_pose} are separate; an orbit event must never call pointer dispatch.

\subsection{Changing the model without changing hardware behavior}
Edit only geometry/material/attachment transforms for visual changes. Keep
attachment objects on \code{carrier} unless the physical part truly belongs to
the fixed stand or azimuth-only support. Check front/side/top/perspective views,
full orbit and zoom, feed/boresight collinearity, and the compass tests.
Hardware behavior belongs in the controller/protocol layers and requires
separate regression evidence.

\section{Virtual pointer transport and trajectory geometry}
\code{VirtualPointer} satisfies the pointer worker's command interface without
creating a serial object or worker thread. It accepts the same 11-byte packets,
validates framing/checksum/opcode, queues one command and reports dispatch to the
controller on the next tick. Absolute targets, four 5-degree jogs and ZERO update
the simulated target. The model normalizes azimuth, bounds elevation to
$[-90,90]$ degrees and advances along the shortest azimuth route with bounded
angular speeds. Virtual bytes use the separate \code{virtual_pointer_tx} capture
channel; command events include a \code{VIRTUAL} source.

\code{VirtualFlight} owns its reference, independent clock, speed and antenna
origin. For a launch-frame ENU point $p_l$, ECEF origins $o_l,o_a$ and ECEF-to-ENU
rotation matrices $R_l,R_a$, the antenna-frame vector is
\[
p_a=R_a(o_l-o_a)+R_a R_l^{\mathsf T}p_l.
\]
Azimuth is $\operatorname{atan2}(E,N)$ modulo $360$ degrees; elevation is
$\operatorname{atan2}(U,\sqrt{E^2+N^2})$. The trajectory's stored origin is
authoritative, and both origins use ellipsoid height. Linear interpolation
uses the existing trajectory contract. Coincident positions have no direction;
exact zenith retains the preceding azimuth.

\code{Controller} snapshots reference identity and antenna coordinates for each
rehearsal. A mismatch stops motion and clears the old rehearsal. The follow path
requires an actual \code{VirtualPointer} instance and cannot dispatch through a
physical worker. Manual actions cancel virtual playback; source changes close
both connection types. Rehearsal never populates the telemetry history or track,
and does not satisfy the ground-station readiness predicate.

\code{widgets.ComboBox} measures polished popup rows before opening, accounting
for frame, padding and scrollbar space while bounding width to the screen.
The stylesheet selects a list popup instead of the macOS menu-style checkmark
delegate. This avoids the clipped mode/camera labels observed with Lucida Grande.
The shared widget serves Settings, mission location and workspace selectors.

\section{Settings storage and application}
\code{AppSettings} is a frozen dataclass, independent of \code{Mission}.
Schema version 1 contains \code{openrocket_jar} (empty means bundled),
\code{sessions_directory} (empty means the per-user sessions directory),
\code{daylight} (false by default), and \code{simulation_timeout} (120 s,
validated integer range 10--1800). Nonempty stored paths must be absolute.
The dialog expands user paths before storage. JAR validation checks the ZIP
container for engine and dependency classes; API compatibility is established
by bridge compilation and actual simulation, not by the file chooser alone.

\code{Controller} loads \code{settings.json} from its data directory.
Malformed settings leave defaults active, preserve the original file and produce
an operator-visible diagnostic. Missing external files remain selected until
repaired; simulation does not silently substitute the bundled engine.
Saving writes and flushes a temporary file in the same directory and then
atomically replaces the settings file. In-memory preferences change only after
a successful write.

\code{MainWindow} installs one application-wide Command-comma QAction with the
native preferences menu role. \code{SettingsDialog} is modeless and reused while
open. Save updates the controller and theme; Cancel has no persistence side
effects. Every simulation receives the immutable settings snapshot at submission.
The private Java runtime and bridge remain bundled; the first classpath entry
is the selected engine. Result metadata records its path, SHA-256 and time limit.
Logging resolves the configured destination when starting a new session.
Portable-flight loading never rewrites machine settings.

\section{Developer setup and native packaging}\label{sec:build}
\subsection{Source environment}
Run the following from the application root, not this report directory:
\begin{lstlisting}[language=bash]
make setup PYTHON=python3.12
make run
make demo
make test
make lint
\end{lstlisting}
Setup creates \code{.venv}, installs pinned build requirements and installs the
package editable without resolving a second dependency set. With uv, an existing
workflow can use \code{uv sync --locked}. Runtime simulation additionally needs
the staged private Java/engine/bridge.

The shared Makefile uses \code{.venv/bin/python} on Unix-like hosts and
\code{.venv/Scripts/python.exe} on Windows. The default target prints help.
\code{make clean} removes generated app \code{build}/\code{dist}; it is not
a session-data purge command. Preserve validation logs you want to retain before
cleaning the application build directory.

\subsection{Runtime and bridge}
\begin{lstlisting}[language=bash]
make runtime
make bridge OPENROCKET_JAR="/absolute/path/to/OpenRocket-MIT-v6.2.jar"
# Or resolve the latest stable release from the team repository:
make bridge OPENROCKET_DOWNLOAD=latest
\end{lstlisting}
\code{runtime} selects OS and process architecture, downloads a Temurin Java 17
JRE, verifies the publisher-provided SHA-256 and stages \code{vendor/java}.
Its manifest records the version, source, hash, OS and architecture. Reuse is
allowed only for a matching target. Stage dependencies/runtime in advance for
offline builds; do not copy a Mac virtual environment or JRE to Windows.

\code{bridge} needs JDK 17+ tools \code{javac} and \code{jar}, and the team's
complete fork JAR. It compiles with Java 17 release and UTF-8 settings, builds
\code{rocket-bridge.jar}, records engine/bridge hashes and copies bridge source
into the vendor tree. With a previously staged engine, later \code{make bridge}
can reuse it. Rebuild the bridge after editing its Java source.

\subsection{Release downloads and engine provenance}
The team's \href{https://github.com/MDNich/ActiveControl_MIT_RktTeam/releases}{GitHub releases} provide the bundled engine. Set
\code{OPENROCKET_DOWNLOAD=latest} to resolve the latest stable release.
\code{OPENROCKET_DOWNLOAD=v6.2} pins a release tag. Asset selection requires the
exact name \code{OpenRocket-MIT-vX.Y.jar}; installer and disk-image assets are
ignored. Supplying both download and local-JAR options is an error.
The downloader checks asset size and any published SHA-256; an absent published
digest is explicitly reported and the downloaded hash is still recorded.
The v6.2 JAR has a verified published checksum.

Preparation stages the candidate and compiles the bridge in a temporary
directory. Download, checksum or compilation failure preserves the previously
staged engine and bridge. The engine manifest records the release tag, asset
URL, source archive URL, published digest, verification status and engine/bridge
hashes. Latest is resolved only on an explicit build request; neither app
startup nor Settings performs automatic updates. A future incompatible release
can fail bridge compilation; production builds should pin a tested tag.

\begin{lstlisting}[language=bash]
make package-macos OPENROCKET_DOWNLOAD=latest
make package-macos OPENROCKET_DOWNLOAD=v6.2
make package-macos OPENROCKET_JAR="/path/to/OpenRocket-MIT-v6.2.jar"
\end{lstlisting}
These options also apply to \code{package-windows}, \code{package} and native
\code{all} targets. With an engine option, packaging stages the engine, rebuilds
the bridge and prepares the private Java runtime before freezing. With neither
option, existing vendor files are reused without a download. The current Mac
package embeds the v6.2 release; release provenance does not itself establish a
reproducible source-to-binary build.

\subsection{Native platform builds}
\begin{lstlisting}[language=bash]
# On macOS, after setup:
make -f Makefile.macos all \
  OPENROCKET_JAR="/absolute/path/to/swing-24.12.RC.01-all.jar"

# On Windows with GNU Make and x64 Python:
make -f Makefile.windows all \
  OPENROCKET_JAR="C:/OpenRocket/swing-24.12.RC.01-all.jar"
\end{lstlisting}
Each platform's \code{all} target runs tests, stages runtime, builds bridge,
packages and verifies the package. PyInstaller builds on its target OS; this is
not cross-compilation. Windows ARM hosts must use x64 Python for the supported
emulated build path; runtime architecture follows the Python process rather
than host hardware alone.

Without GNU Make, Windows PowerShell can call the scripts directly:
\begin{lstlisting}
py -3.12 scripts/build.py setup
.venv\Scripts\python.exe -m pytest -q
.venv\Scripts\python.exe scripts/build.py runtime
.venv\Scripts\python.exe scripts/build.py bridge --engine "C:\OpenRocket\swing-all.jar"
.venv\Scripts\python.exe scripts/build.py package --platform Windows
.venv\Scripts\python.exe scripts/verify_package.py
\end{lstlisting}
The prior Windows VM used Python 3.13 x64; use the installed matching launcher
version. This historical validation does not assert that the newest Mac feature
revision has been rebuilt there.

\subsection{Package contents and artifacts}
\src{packaging/RocketGNCMonitor.spec} includes the resources, in-app docs and vendor
tree, FFmpeg data, package metadata and the native MGRS library. It excludes
unneeded Qt WebEngine/QML/Quick and scientific packages. The launcher is
windowed. macOS bundling declares camera usage and the application identity.

The build renders the repository SVG icon into native icons, gathers dependency
license/notice files, writes \code{vendor/build-manifest.json} and retains
FFmpeg build information. macOS produces
\nolinkurl{dist/RocketGNCMonitor.app} and an architecture-labeled DMG. Windows produces
a complete \code{dist/RocketGNCMonitor} folder and ZIP. Distributable archives
have SHA-256 sidecars.

\code{MAC_SIGN_IDENTITY} can select a team Developer ID identity at build time;
the current local build uses ad-hoc signing. Notarization/stapling and Windows
release signing are external release steps. Preserve existing third-party
notices, the proprietary font provenance, Java legal files, engine and bridge
source/provenance materials. The exact corresponding engine source remains an
explicit release record to resolve; this documentation does not grant new
licenses or credentials.

\section{Testing, verification and evidence}\label{sec:tests}
\subsection{Test organization}
\begin{longtable}{@{}L{.32\linewidth}L{.61\linewidth}@{}}
\toprule Test group & Contract exercised \\\midrule\endhead
\code{test_protocol.py}, \code{test_serial.py} & Decoder corruption/fragmentation, clock and sequence behavior, pointer bytes and OS pseudo-terminal transport. \\
\code{test_parity.py} & Captured old-UI commands/CSV/shortcuts, actual Qt key events, independent ports and retained confirmations. \\
\code{test_controller.py} & Startup locks, source isolation, stale tracking, dispatch/pending state and mode behavior. \\
\code{test_domain.py}, \code{test_location.py} & Geographic/datum/wind validation, MGRS/Plus Codes, preset and coordinate persistence. \\
\code{test_demo.py}, \code{test_legacy.py} & Exact bundled logs, timing/seek/quality behavior, safe CSV conversion and provenance. \\
\code{test_recording.py}, \code{test_session_flow.py} & SQLite/raw/CSV recovery, pending-batch drain, record-stop-replay and video seek. \\
\code{test_dual_video.py} & Camera reservations and dropdown exclusion, independent streams/failures, legacy replay, dual-video archive round trip and minimum-size panels. \\
\code{test_flight.py} & Portable assets, active snapshots, full-demo and buffer exports, corruption/path rejection, atomic failure and final video inclusion. \\
\code{test_media.py} & FFmpeg arguments, capture/recording lifecycle and available runtime checks. \\
\code{test_ui.py} & Workspace rendering, legacy values/readiness, orbit independence, compass chirality and pseudo-terminal radio status. \\
\code{test_simulation.py} & Bundled N8406 identity, real engine trajectory, named missing-motor and unpowered errors. Engine-dependent cases skip if runtime/bridge is absent. \\
\code{test_virtual_pointer.py} & Packet actuation, frame geometry, playback, manual takeover, physical isolation, profile/flight persistence and Qt controls. \\
\code{test_settings.py} & Shortcut, persistence, cancel/defaults, invalid settings, recording destination and per-job snapshots. \\
\code{test_build.py} & Runtime architecture; latest/pinned release selection, checksums, local JARs and failure preservation. \\
\bottomrule
\end{longtable}

Run focused tests for the seam being changed, then the appropriate full release
checks. Offscreen Mac/Linux Qt tests can use:
\begin{lstlisting}[language=bash]
QT_QPA_PLATFORM=offscreen .venv/bin/python -m pytest -q
.venv/bin/python -m ruff check src tests --select F
git diff --check
\end{lstlisting}
Windows tests should exercise the native Qt Windows platform; its offscreen
font behavior is not equivalent to checking a real native window.
These serial tests use pseudo-terminals or in-memory substitutes, not physical
boards or recovery hardware.

\subsection{Legacy parity baseline}
\code{tests/fixtures/legacy_reference.json} was captured from the original
\code{UI.py}, \code{rocket.py} and \code{pointer.py}. It records source hashes and
118 rocket command cases, eight pointer cases, geographic tracking cases,
all 43 CSV columns and the ten original shortcut strings.
The fixture generator strips physical connection behavior and substitutes
in-memory streams. Do not regenerate it merely to make a failing test accept
a changed protocol; first establish the intended compatibility decision.

The retained compatibility details include receiver-height division, five-decimal
frozen GPS, ground height zero for tracking, servo degree conversion/rounding and
two unused padded pyro list entries in good-packet CSV. Documented fixes include
rejected-packet counter accuracy, bad-packet-only field inconsistencies, eight
BMS columns and duplicate old shortcut connections.

\subsection{Packaged verification}
\begin{lstlisting}[language=bash]
make verify-package
.venv/bin/python scripts/verify_package.py \
  --model /absolute/path/to/zephy_testlaunch.ork \
  --output build/package-verification
\end{lstlisting}
The script launches the actual packaged executable with a minimal system PATH
and no developer Python/Java overrides. It checks locked LIVE startup,
Lucida Grande, coordinate conversion, planning/recorded flight round trips with
URRG, all three demo logs plus video, and the bundled example OpenRocket run.
The optional model flag adds a supplied model to the same packaged worker test.

CLI smoke entry points are:
\begin{longtable}{@{}L{.30\linewidth}L{.63\linewidth}@{}}
\toprule Flag & Use \\\midrule\endhead
\code{--demo} & Explicit recorded demo startup. \\
\code{--demo-station GS1|GS2|GS3} & Receiver selection; default GS2. \\
\code{--data-dir PATH} & Override per-user application data directory. \\
\code{--virtual-pointer-smoke CSV} & Exercise virtual connection, actuation, reference playback/seek and disconnect without physical transport. \\
\code{--startup-smoke} & Check disconnected/locked LIVE startup and exit. \\
\code{--smoke-test} & Exercise recorded demo and synthetic video, write report, exit. \\
\code{--flight-smoke} & Save/load planning and GS2 flights, URRG/assets/reference, paused isolated replay. \\
\code{--simulation-smoke MODEL} & Exercise private engine with synthetic zero-coordinate site and write report. \\
\code{--screenshot PATH} & Optional native render for smoke verification; not used for this report's diagrams. \\
\bottomrule
\end{longtable}

\subsection{Recorded results and their limits}
The current virtual-pointer revision passed 101 Mac regression tests. Native Cocoa popup checks cover both themes and long device names. The previous Settings/engine revision passed 94 tests. Settings persistence, shortcut handling, bundled/custom simulations and local/latest/pinned build paths are covered. The earlier dual-video revision passed 77 tests. New checks
cover camera reservations during open/close, reciprocal camera and serial
dropdown exclusion, independent stream failure and frames, old single-video
replay, two-feed flight archives and compact layouts. The refreshed package
smoke report records Digital and Analog frame counts separately.

The retained September 20 Mac validation reports 71 passing regression tests,
followed by focused final checks. Packaged flight verification restored all
5,464 GS2 samples, URRG, a copied model and reference in paused REPLAY with no
physical worker. The bundled example simulation returned 1,551 rows.
The supplied Zephyrus ORK returned 11,626 rows in the package's synthetic
zero-coordinate smoke setup; this is a dependency/configuration test rather
than the actual launch prediction.

The prior Windows 11 ARM VM exercised an x64 package under emulation with
30 tests passing and one POSIX-only skip, before the current revision.
Mac signature and DMG checks establish package integrity, not notarization or
hardware qualification. The copied evidence in this report folder preserves
the original result rather than relabeling it as a fresh test run.

\section{Extension recipes and maintenance boundaries}\label{sec:extend}
\subsection{Add another launch-site preset}
Add the name/code to \code{LaunchLocation}'s preset selector and its selection
handler, use the existing converter, emit \code{preset_selected} and preserve
canonical coordinates. Keep elevations explicit. Add a test for exact input
provenance, save/load stability, display-format changes, manual edit clearing
and cancellation. If presets grow, extracting a data-only catalog is a natural
refactor; do not introduce a required network geocoder for offline presets.

\subsection{Support the new three-axis telemetry board}
Define a versioned wire contract first: framing, checksum/CRC, units, axes,
timestamps, reset behavior, sequence numbering and channel identities. Implement
a separate decoder/adapter that produces the established \code{Sample} contract.
Populate measured/demand fields only when the firmware supplies them.
Do not reinterpret legacy servo channel numbers as canards/tabs without an
explicit hardware mapping.

Add golden packet vectors, fragmented/corrupt/reboot tests, recording and archive
round trips, replay/UI tests and captured new-board data. Preserve the old
Zephyrus adapter so existing logs and hardware retain their behavior. Update the
protocol documentation and qualification status with the code.

\subsection{Add a command or pointer feedback protocol}
Keep pure encoding in the protocol/command modules and dispatch authorization in
the controller. Add the appropriate UI control and confirmation behavior,
reference bytes and no-duplicate-shortcut tests. Define acknowledgment IDs,
timeouts and partial-write semantics explicitly before showing a new
\emph{confirmed} status.

For measured pointer feedback, create a structured RX parser and a separate
measured-pose field. Do not overwrite the meaning of the last-sent target.
Mechanical homing, cable limits and motion stop require an agreed firmware and
operational contract, not merely a renamed host button.

\subsection{Extend a saved schema}
For an optional mission field, add a dataclass default, validation, editor if
needed, and mission/flight/session tests. Older version-1 files rely on default
values for absent known fields. Breaking semantic changes need a version bump
and explicit migration or rejection; filtering unknown fields is not migration.

For new sample shapes or archive files, update \code{Sample}, import adapters,
record/export paths, \code{flight.load_flight} validation and replay consumers.
Exercise save during active logging, failure before replacement and reopening
without original source assets. Preserve queue bounds and authoritative capture
independence from UI history.

\subsection{Change simulation physics}
Modify the Java bridge only against the matching fork API and rebuild
\code{rocket-bridge.jar}. Record solver/options, tests and model/motor hashes.
Controlled canard/tab dynamics need explicit model parameters and physics
qualification; do not simply clear the \code{controlled_model_validated:false}
flag or infer missing vehicle feedback from demo traces.

\subsection{A practical change checklist}
Before distributing a change, review the affected source-mode behavior and
resource ownership; run meaningful focused tests; retain legacy parity where
applicable; inspect affected windows at 1120 by 800 and larger sizes; rebuild the
native package; run packaged verification; and update operator/developer docs
and release evidence. Physical interface or control changes additionally require
the team's hardware acceptance process. Never use a real recovery command as a
software smoke test.

\appendix
\section{Complete mission configuration reference}\label{app:mission}
The table follows the current \code{Mission} dataclass, including retained
compatibility fields. Mission JSON adds \code{schema_version:1}; the default
values below are application defaults, not a recommended launch configuration.
The validator applies finite-number/range checks where defined; UI spinbox
limits and dataclass validation are not identical for every field.
\begin{longtable}{@{}L{.28\linewidth}L{.18\linewidth}L{.47\linewidth}@{}}
\toprule Field & Default & Meaning \\\midrule\endhead
\code{name} & Untitled mission & Human-readable mission name. \\
\code{latitude} & 0.0 & WGS84 launch latitude in degrees; -90 to 90. \\
\code{longitude} & 0.0 & WGS84 launch longitude in degrees; -180 to 180. \\
\code{altitude} & 0.0 & Launch ellipsoid altitude in metres; used by geographic transforms. \\
\code{altitude_msl} & 0.0 & Launch mean-sea-level elevation in metres; weather and engine atmosphere. \\
\code{site_configured} & False & Explicit origin-established flag; required for LIVE trajectory and simulation. \\
\code{launch_location_format} & latlon & Entry provenance: latlon, mgrs or pluscode. \\
\code{launch_location_code} & (empty) & Entered/resolved code; up to 80 characters. Canonical lat/lon remains authoritative. \\
\code{launch_site_name} & (empty) & Preset name, currently URRG or empty for Custom; up to 80 characters. \\
\code{legacy_altitude} & unknown & unknown, ellipsoid, gps\_agl or barometric\_agl; selects trajectory height interpretation. \\
\code{canard_count} & 4 & Number of UI canard slots, integer 0 to 16; four tabs are always added. \\
\code{pointer_latitude} & 0.0 & Virtual antenna WGS84 latitude, degrees; physical tracking retains frozen receiver GPS. \\
\code{pointer_longitude} & 0.0 & Virtual antenna WGS84 longitude, degrees. \\
\code{pointer_altitude} & 0.0 & Virtual antenna WGS84 ellipsoid altitude, metres; physical legacy tracking retains height zero. \\
\code{pointer_site_configured} & False & Boolean assertion that the antenna location is established; required for virtual trajectory tracking. \\
\code{pointer_calibrated} & False & Retained assertion, always cleared by disk load and mode changes; not a manual-send prerequisite. \\
\code{pointer_az_offset} & 0.0 & Retained offset field; not applied by current legacy manual/tracking dispatch. \\
\code{pointer_el_offset} & 0.0 & Retained offset field; not applied by current legacy manual/tracking dispatch. \\
\code{pointer_el_min} & 0.0 & Stored route-helper envelope minimum; not an enforced operator control limit. \\
\code{pointer_el_max} & 80.0 & Stored route-helper envelope maximum; ordered within -90 to 90 degrees. \\
\code{pointer_az_min} & 0.0 & Stored route-helper azimuth minimum; not enforced by the legacy dispatch path. \\
\code{pointer_az_max} & 359.0 & Stored route-helper azimuth maximum; ordered within 0 to 360 degrees. \\
\code{pointer_full_rotation} & False & Stored route-helper choice; not a firmware cable-wrap guarantee. \\
\code{freshness} & 2.0 & Maximum acceptable sample age, 0.1 to 60 seconds. \\
\code{video_offset} & 0.0 & Seconds; positive value delays recorded video relative to telemetry. \\
\code{wind} & One calm layer at height 0 & 1 to 200 increasing nonnegative AGL heights with east/north velocity components; max speed 150 m/s. \\
\code{wind_source} & Manual · calm & Human-readable profile provenance. \\
\code{weather_raw} & \{\} & Cached normalized request/result/provider payload, when retrieved. \\
\code{model} & (empty) & Selected ORK path. Archives embed it and rebind the extracted path. \\
\code{motor_files} & [] & Up to 100 custom curve paths; selected files are embedded in flight archives. \\
\code{simulation_index} & 0 & Saved ORK simulation, zero-based integer 0 to 1000. \\
\code{rail_length} & 2.0 & Launch rod length in metres; greater than 0 and at most 100. \\
\code{rail_tilt} & 0.0 & Tilt from vertical in degrees; at least 0 and less than 90. \\
\code{rail_heading} & 0.0 & True-north bearing in degrees; UI range 0 to 359.99. \\
\code{seed} & 42 & Integer random seed, 0 to 2147483647. \\
\code{low_battery} & 9.0 & Total-battery alert threshold in volts; 1-second dwell and 0.3 V clearing hysteresis. \\
\bottomrule\end{longtable}


\section{Complete normalized sample reference}\label{app:sample}
\code{Sample.from_dict} retains known dataclass fields and uses defaults for
absent optional values. The file formats add their own structural validation;
the dataclass constructor alone is not a complete untrusted-input validator.
Vector and dictionary meanings are defined by the producing protocol adapter.
\begin{longtable}{@{}L{.22\linewidth}L{.20\linewidth}L{.51\linewidth}@{}}
\toprule Field & Default & Meaning \\\midrule\endhead
\code{t} & required & Device uptime seconds after wrap extension, or imported equivalent; align separately to flight zero. \\
\code{sequence} & required & Packet counter; LIVE wire counter is uint16. \\
\code{source} & required & Provenance such as LIVE, DEMO or LEGACY\_CSV. REPLAY is a controller mode, not a rewrite of sample source. \\
\code{received} & runtime factory & Host monotonic receive time in seconds; replay refreshes this for current display freshness. \\
\code{utc} & runtime factory & Unix reception timestamp in seconds; original UTC retained for imported/demo data. \\
\code{altitude} & None & Reported filtered barometric altitude, metres, or unavailable. \\
\code{velocity} & None & Reported integrated velocity, metres/second, or unavailable. \\
\code{latitude} & None & Rocket WGS84 latitude in degrees, or unavailable. \\
\code{longitude} & None & Rocket WGS84 longitude in degrees, or unavailable. \\
\code{gps_altitude} & None & Reported GPS height in metres, with explicit mission interpretation. \\
\code{gps_fix} & 0 & GPS fix code; geographic display requires at least 3. \\
\code{attitude} & None & Optional three-element roll/pitch/yaw integrated rotations, degrees; not a qualified attitude solution. \\
\code{rates} & None & Optional XYZ angular rates, degrees/second; legacy Y sign is retained. \\
\code{acceleration} & None & Optional XYZ acceleration, metres/second squared. \\
\code{enu} & None & Optional local east/north/up metres; frame/provenance must be known. \\
\code{battery} & None & Optional total battery volts, sum of the three cells for LIVE. \\
\code{rssi} & None & Optional receiver RSSI in dBm. \\
\code{phase} & Unknown & Received or imported state label; unknown codes remain visible. \\
\code{actuators} & runtime factory & Dictionary by channel name; demand/measured angles or legacy drive when actually available. \\
\code{details} & runtime factory & Extensible decoded legacy values, boot/time metadata, receiver data, raw CSV provenance and quality notes. \\
\bottomrule\end{longtable}


\section{Legacy CSV columns and derived values}\label{app:csv}
This is the exact 43-column order written by the compatibility adapter.
The rich \code{SessionReader.export_csv} format is different and includes full
sample JSON. Imported records retain their original row text and available
decoded values. List padding described below refers to live good-packet output;
it should not be used to invent extra hardware channels.
\begin{longtable}{@{}L{.08\linewidth}L{.34\linewidth}L{.51\linewidth}@{}}
\toprule Column & Field & Meaning \\\midrule\endhead
1 & \code{timestamp} & Receive UTC, Unix seconds. \\
2 & \code{pyros} & Six pyro continuity codes. \\
3 & \code{servos} & Four legacy drive values. \\
4 & \code{servos_deg} & Four converted servo angles; first two use 60-degree scale, last two 50-degree scale; good CSV rounded to two decimals. \\
5 & \code{accelerometer} & XYZ acceleration, m/s squared. \\
6 & \code{barofilteredalt} & Filtered barometric altitude, m. \\
7 & \code{temp} & Legacy calibrated temperature, degrees C. \\
8 & \code{gyro} & XYZ gyro rates, degrees/s. \\
9 & \code{gps_fix} & Rocket fix code. \\
10 & \code{lat} & Rocket latitude, degrees. \\
11 & \code{lon} & Rocket longitude, degrees. \\
12 & \code{gpsalt} & GPS height, m; datum depends on firmware/configuration. \\
13 & \code{gps_horiz_prec} & Horizontal GPS accuracy, m. \\
14 & \code{gps_vert_prec} & Vertical GPS accuracy, m. \\
15 & \code{gps_num_sat} & Satellite count. \\
16 & \code{flight_time} & Original raw device clock, milliseconds. \\
17 & \code{yaw_gyro_int} & Integrated yaw, degrees. \\
18 & \code{pitch_gyro_int} & Integrated pitch, degrees. \\
19 & \code{roll_gyro_int} & Integrated roll, degrees. \\
20 & \code{state} & Legacy state enum string. \\
21 & \code{pktnum} & Packet counter. \\
22 & \code{rssi} & Ground receiver RSSI, dBm. \\
23 & \code{armed_pyros} & Armed bits as list; good-packet CSV retains two unused zero entries. \\
24 & \code{fired_pyros} & Fired bits as list; same good-packet padding. \\
25 & \code{badpackets} & Rejected-packet counter. \\
26 & \code{rxrssi} & Onboard RSSI, dBm. \\
27 & \code{accel_integrated_velo} & Reported integrated velocity, m/s. \\
28 & \code{baro_max_alt} & Maximum barometric altitude, m. \\
29 & \code{gps_max_alt} & Maximum GPS altitude, m. \\
30 & \code{pyro_resistances} & Pyro resistance list; good-packet output pads two unused zero entries. \\
31 & \code{cell_voltages} & Three cell voltages, V. \\
32 & \code{total_current} & Total current, A. \\
33 & \code{converter_voltages} & Six rail voltages, V. \\
34 & \code{converter_currents} & Six rail currents, A. \\
35 & \code{bms_protections_enabled} & BMS enabled protection bits. \\
36 & \code{bms_protection_status} & BMS triggered protection bits. \\
37 & \code{bms_temp} & BMS temperature, degrees C. \\
38 & \code{enabled_status} & Derived rail-enabled states from voltage proximity to nominal. \\
39 & \code{angleFromVertical} & Legacy angle-from-vertical derived from pitch/yaw, degrees. \\
40 & \code{gnd_lat} & Frozen/display receiver latitude convention, rounded to five decimals. \\
41 & \code{gnd_lon} & Receiver longitude convention, rounded to five decimals. \\
42 & \code{gnd_fix} & Legacy ground fix Boolean: trailer fix equals 3. \\
43 & \code{gnd_alt} & Legacy receiver height: signed wire integer divided by 1000, rounded for display/CSV. \\
\bottomrule\end{longtable}


The live servo degree conversion is
\[
 \alpha_i=\frac{d_i-1500}{500}
 \begin{cases}60^\circ,& i=0,1,\\50^\circ,& i=2,3.\end{cases}
\]
The legacy angle from vertical is
\(\arccos(\cos p\cos y)\), with the cosine product clamped to \([-1,1]\).
The raw temperature calibration in \code{legacy_values} is
\[
 T=\frac{2000+(T_{\mathrm{raw}}-\texttt{0x91E3}\cdot256)
                    \texttt{0x6FEC}/2^{23}}{100}.
\]
These formulas reproduce the former monitor; they do not establish new sensor
calibration accuracy.

\section{Worked file examples}\label{app:examples}
The \code{examples} directory accompanies the source report. Its inputs have been
loaded through the application's validators as documentation examples.

\subsection{URRG planning template}
\code{examples/urrg-mission.json} includes all current mission fields and the
exact URRG code/canonical coordinates. It deliberately leaves
\code{site_configured:false}, altitude convention unknown and both elevations
zero so the example cannot be mistaken for a surveyed setup. Open it with
\ui{Open mission}, then review and complete those values. The essential location
portion is:
\begin{lstlisting}
{
  "schema_version": 1,
  "name": "URRG template - verify launch elevations",
  "launch_site_name": "URRG",
  "launch_location_format": "mgrs",
  "launch_location_code": "18TUN2061530290",
  "site_configured": false,
  "legacy_altitude": "unknown"
}
\end{lstlisting}
The full file also supplies the decoded latitude and longitude; do not treat
this shortened excerpt as a replacement for the full template.

\subsection{Wind profile}
\lstinputlisting{examples/wind.json}
The second layer has eastward and southward components. This is a formatting
example, not a retrieved URRG forecast.

\subsection{Normalized trajectory}
\lstinputlisting{examples/reference.csv}
The matching manifest is:
\lstinputlisting{examples/reference.json}
The example is explicitly synthetic and is therefore rejected as a LIVE
reference. It illustrates file structure and can be used for offline tests.

\subsection{Portable flight metadata excerpt}
\begin{lstlisting}
{
  "format": "rocket-gnc-flight",
  "schema_version": 1,
  "app_version": "0.1.0",
  "source_mode": "DEMO",
  "scope": "Complete Zephyrus GS2 dataset",
  "session": "session",
  "samples": 5464,
  "position": 174.0,
  "flight_zero": 747.766,
  "time_aligned": true,
  "files": {
    "mission.json": {
      "size": "<actual byte count>",
      "sha256": "<actual content hash>"
    }
  }
}
\end{lstlisting}
This explanatory excerpt omits other entries and contains descriptive
placeholders; it is not an importable archive. Use the app's Save flight
operation to create a valid inventory and package.

\section{Source and validation index}\label{app:evidence}
The code is the behavioral authority for this report. The following retained
documents supplement the source map and explain provenance and compatibility:
\begin{longtable}{@{}L{.35\linewidth}L{.58\linewidth}@{}}
\toprule Source & Role \\\midrule\endhead
\src{docs/OPERATOR_GUIDE.md} & In-app operator help; this report expands it into full functionality and developer coverage. \\
\src{docs/LEGACY_PARITY.md} & Original UI inventory, captured vectors, exact shortcuts and compatibility exceptions. \\
\src{docs/PROTOCOL.md} & Audited telemetry, pointer and uplink wire contract. \\
\src{docs/DEMO_DATA.md} & GS1/GS2/GS3 source identity, timing and known quality issues. \\
\src{docs/BUILD.md} & Native build, runtime staging, offline preparation and signing workflow. \\
\src{docs/IMPLEMENTATION_STATUS.md} & Current versus historical acceptance status. \\
\src{docs/THIRD_PARTY_NOTICES.md} & Bundled dependency notices and provenance references. \\
\src{resources/fonts/README.md} & User-designated font resource and unchanged source hash. \\
\src{resources/demo/manifest.json} & Original CSV hashes, row counts, durations and launch alignment. \\
\src{resources/motors/manifest.json} & Bundled motor resource integrity and source identity. \\
\src{tests/fixtures/legacy_reference.json} & Captured original protocol/CSV/shortcut behavior. \\
\code{evidence/source-manifest.json} & Exact reviewed source/resource bytes and report-style reference hash. \\
\nolinkurl{evidence/application-validation-settings-2026-09-21.json} & Current 94-test Settings/engine validation and packaged Mac checks. Earlier evidence files retain their historical context. \\
\bottomrule
\end{longtable}

The original design plans describe intentions and earlier stages; they are not
used here to claim unimplemented capabilities. In particular, active legacy
manual pointing does not enforce the early proposed software envelope, current
DEMO uses recorded data rather than the synthetic fixture, and current Windows
acceptance is older than this Mac revision.

\section{Glossary}
\begin{longtable}{@{}L{.23\linewidth}L{.70\linewidth}@{}}
\toprule Term & Meaning in this app \\\midrule\endhead
AGL & Height above the configured launch/ground reference. \\
Ellipsoid height & Height relative to the WGS84 ellipsoid; used in ECEF/ENU transforms. \\
MSL & Mean sea level elevation; entered separately for atmosphere and weather. \\
ECEF / ENU & Earth-centered Earth-fixed / local east-north-up coordinate systems. \\
GNC & Guidance, navigation and control; this app monitors available data and sends legacy commands. \\
MGRS / Plus Code & Offline geographic coordinate representations; neither supplies elevation. \\
Mission & Configuration object; mission JSON alone can contain external asset paths. \\
Session & Directory containing indexed recorded data, raw bytes, CSVs, events and optional video/reference. \\
Flight file & Portable versioned ZIP64 archive with mission/assets and available recording/reference data. \\
Snapshot & Consistent committed subset saved while acquisition continues; active video segment may be absent. \\
Nominal reference & OpenRocket output under the recorded options; not qualified controlled-vehicle dynamics. \\
Sent / acknowledged & Serial dispatch / confirmed device response. The legacy command paths provide the first, not the second. \\
Source mode / phase & LIVE, DEMO or REPLAY / received onboard rocket state; independent concepts. \\
WAL & SQLite write-ahead log; retain it with a crashed source session during recovery. \\
\bottomrule
\end{longtable}

\label{LastContentPage}
\end{document}
